\documentclass[11pt]{article}
\usepackage[margin=1in]{geometry}
\usepackage{amsmath,amssymb,amsthm,mathtools}
\usepackage{hyperref}
\usepackage{xcolor}
\hypersetup{colorlinks=true,linkcolor=blue,citecolor=blue,urlcolor=blue}
\setlength{\emergencystretch}{2em}

% ---------------------------------------------------------------------------
% Blueprint macros.  When this file is later compiled with the `leanblueprint`
% tool (https://github.com/PatrickMassot/leanblueprint) these macros are
% provided by the tool and drive the dependency graph.  Until then they are
% plain-text annotations so the document compiles with plain pdflatex.
% ---------------------------------------------------------------------------
\providecommand{\lean}[1]{\leavevmode\par\nopagebreak
  {\noindent\normalfont\footnotesize\raggedright\texttt{Lean: #1}\par}\nopagebreak}
\providecommand{\leanok}{\par\nopagebreak
  {\noindent\normalfont\footnotesize\textcolor{green!50!black}{[formalised]}\par}\nopagebreak}
\providecommand{\uses}[1]{\par\nopagebreak
  {\noindent\normalfont\footnotesize Uses: #1\par}\nopagebreak}
\providecommand{\notready}{}

\newtheorem{theorem}{Theorem}[section]
\newtheorem{proposition}[theorem]{Proposition}
\newtheorem{lemma}[theorem]{Lemma}
\newtheorem{definition}[theorem]{Definition}
\newtheorem{remark}[theorem]{Remark}

\newcommand{\E}{\mathbb{E}}
\newcommand{\R}{\mathbb{R}}
\newcommand{\N}{\mathbb{N}}
\newcommand{\Sig}{\Sigma_N}
\newcommand{\bra}[1]{\langle #1\rangle}
\DeclareMathOperator{\cov}{Cov}

\title{The Parisi formula in Lean 4\\ \large Blueprint (working draft)}
\author{ParisiFormula project}
\date{September 2026}

\begin{document}
\maketitle

\section*{Purpose and status}

This document is the mathematical plan for a Lean~4 formalisation of the Parisi
formula for the Sherrington--Kirkpatrick (SK) model, following Guerra's interpolation
and Talagrand's proof \cite{Tal06}.  Each statement is paired with the name of its
Lean counterpart where one exists. Statements marked \emph{[formalised]} have checked
proofs under their stated hypotheses; open targets and conditional deductions are
identified explicitly. The active theorem statement is in
\texttt{Targets/Talagrand.lean}; its supporting modules are described in the roadmap.
\texttt{Targets/Milestones.lean} also contains older, off-path targets.

\textbf{Status (7 September 2026).} The exact-covariance SK version of Theorem~2.1,
Guerra's RSB upper bound, and Theorem~2.4 are proved. A downstream integration
theorem combines Theorem~2.4, Proposition~2.3 and the checked convergence argument
to prove exactly the mathematical conclusion of Theorem~2.2. All these results
pass the standard-axiom audit in \texttt{Targets/GuerraAudit.lean}. The refactored
final module exports this result as the canonical Theorem~2.2 and uses it to prove
the Parisi formula. Thus the full critical path is checked in the stated
exact-covariance SK setting.
Three older placeholders remain off this critical path: Guerra--Toninelli
monotonicity, the standalone thermodynamic-limit target, and a uniform Lipschitz
estimate for the scheme. Fixed-level continuity and existence of a minimizer are proved.
Stationarity now holds at every level after exact scheme reduction. The actual
endpoint Hessian calculus and depth-uniform regularity are checked. Lemma~4.9's
interior third derivative and closed-interval Taylor control now supply
Proposition~4.10 without an assumed regularity bound. The local-left and initial
estimates (Propositions~5.1 and~5.3) are checked. The local-right estimate
(Proposition~5.2) is checked for every level of a reduced scheme, including
zero first overlap and the terminal interval. The far-right strict improvement
(Proposition~5.6) is also checked, with a deficit chosen before system size
and disorder. The negative initial interval (Proposition~5.4) is now checked
for reduced schemes by exact SK reflection and conditioning, preserving the
original shared field and free energy. For Proposition~5.7, strict scalar
interchange and sorting, the mixed paired comparison, sorted scalar
identification, the original-free-energy endpoint and the time-zero deficit
are checked. The actual mixed-pressure derivative inequality and its transport
to the original free energy are now checked. Proposition~5.7's outside-neighbor
bound is checked for the reduced near-minimizing schemes used in the proof,
including all trial intervals, both signs, zero first overlap and time zero.
Its deficit is independent of system size; uniform assembly into
Theorem~2.4 remains open.

Dependencies are Mathlib and the pinned RSAT library
(\texttt{njimaMath/research\_public}), with local ports originating in
\texttt{or4nge19/SpinGlass} (M.~Cipollina). The root \texttt{Lemmas/} and
\texttt{port/} directories are historical copies, not build targets.
See \texttt{README.md} for the current module map and \texttt{docs/PROVENANCE.md}
for dependency pins and source origins.

\textbf{Milestones.} (1) Existence of the thermodynamic limit (Guerra--Toninelli).
(2) The finite-step Parisi functional and its continuity.
(3) Guerra's replica-symmetry-breaking upper bound.
(4) Talagrand's lower bound, hence the Parisi formula.
These are historical project labels. The active next step is Theorem~2.2 in
Milestone~4; the independent Milestone~1 proof is not a prerequisite of the current
deduction of the Parisi formula.

\section{Conventions}

\begin{definition}[Configurations and overlap]
\lean{SpinGlass.Config, SpinGlass.spin, SpinGlass.overlap}\leanok
For $N\ge1$ let $\Sig=\{-1,1\}^N$. For $\sigma,\tau\in\Sig$ the overlap is
$R(\sigma,\tau)=\frac1N\sum_{i\le N}\sigma_i\tau_i$.
\end{definition}

\begin{definition}[SK disorder]
\lean{SpinGlass.SKDisorder, SpinGlass.sk\_cov\_kernel}\leanok
An SK disorder at inverse temperature $\beta$ is a centred Gaussian family
$(U(\sigma))_{\sigma\in\Sig}$ with
\[
\E\,U(\sigma)U(\tau)=\frac{N\beta^2}{2}\,R(\sigma,\tau)^2 .
\]
Equivalently this is $N\xi(R)$ with $\xi(x)=\beta^2x^2/2$.
The local Theorem~2.1 is specialized to this kernel, not a general mixed $p$-spin model.
The Gaussian structure is encoded through a covariance operator on the Hilbert space
$\R^{\Sig}$ (\texttt{PhysLean.Probability.GaussianIBP.IsGaussianHilbert}).
\end{definition}

\begin{definition}[Free energy]
\lean{SpinGlass.free\_entropy, SpinGlass.free\_energy\_density}\leanok
With external field $h\in\R$,
\[
Z_N=\sum_{\sigma\in\Sig}\exp\Big(U(\sigma)+h\sum_{i\le N}\sigma_i\Big),
\qquad
F_N(\beta,h)=\frac1N\,\E\log Z_N .
\]
\end{definition}

\begin{remark}
The current formal statements assume the exact covariance above. Consequently the
proved interpolation identity and upper bound have no finite-size covariance-error
term. Extending the formalisation to a general covariance or proving a bridge from
another finite-volume normalization is not part of the completed Theorem~2.1 claim.
\end{remark}

\section{Milestone 1: the thermodynamic limit}

The annealed bound, endpoint algebra, covariance comparison, and conditional
Fekete assembly are proved. Proposition~\ref{prop:mono} and the standalone limit
target remain open. This section records an off-path development; it is not a
remaining component of Theorem~2.1.

\begin{lemma}[Annealed bound]\label{lem:annealed}
\lean{SpinGlass.Targets.free\_entropy\_le\_annealed}\leanok
For all $N\ge1$, $F_N(\beta,h)\le\log2+\frac{\beta^2}{4}+|h|$.
\end{lemma}
\begin{proof}
By Jensen, $\E\log Z_N\le\log\E Z_N$.  Since $\E\exp U(\sigma)=\exp(N\beta^2/4)$ and
$|h\sum_i\sigma_i|\le|h|N$, we get $\E Z_N\le 2^N\exp(N\beta^2/4+|h|N)$.
The bound $|\sum_i\sigma_i|\le N$ and the Gaussian moment generating function are
provided by the upstream library and \texttt{ParisiFormula/AnnealedBound.lean}.
\end{proof}

\begin{definition}[Guerra--Toninelli interpolation]
\lean{SpinGlass.K\_block, SpinGlass.K\_interpol, SpinGlass.\ensuremath{\Phi}}\leanok
Let $U_{N+M},U_N,U_M$ be independent SK disorders on $N+M$, $N$, $M$ spins.
Write $\gamma=(\alpha,\sigma)\in\Sigma_{N+M}$ with $\alpha\in\Sig,\sigma\in\Sigma_M$, and
\[
K_t(\gamma)=\sqrt t\,U_{N+M}(\gamma)+\sqrt{1-t}\,\big(U_N(\alpha)+U_M(\sigma)\big),
\qquad
\Phi(t)=\E\log\sum_\gamma\exp\Big(K_t(\gamma)+h\sum_i\gamma_i\Big).
\]
\end{definition}

\begin{lemma}[Endpoints]
\lean{SpinGlass.Z\_interpol\_one, SpinGlass.Z\_interpol\_zero}\leanok
$\Phi(1)=(N+M)\,F_{N+M}$ and $\Phi(0)=N F_N+M F_M$.
\end{lemma}

\begin{lemma}[Covariance comparison]\label{lem:cov}
\lean{SpinGlass.cov\_deriv\_diag, SpinGlass.cov\_deriv\_offdiag\_nonpos}\leanok
Let
\begin{align*}
C_L(\gamma,\gamma')&=\frac{(N+M)\beta^2}2R_{N+M}(\gamma,\gamma')^2,\\
C_{\mathrm{blk}}(\gamma,\gamma')&=\frac{N\beta^2}2R_N(\alpha,\alpha')^2
  +\frac{M\beta^2}2R_M(\sigma,\sigma')^2.
\end{align*}
Then $C_L-C_{\mathrm{blk}}$ vanishes on the diagonal and is $\le0$ off the diagonal.
\end{lemma}
\begin{proof}
On the diagonal all overlaps equal $1$. Off the diagonal, with $\lambda=N/(N+M)$,
$R_{N+M}=\lambda R_N+(1-\lambda)R_M$, and convexity of $x\mapsto x^2$ gives
$R_{N+M}^2\le\lambda R_N^2+(1-\lambda)R_M^2$.
\end{proof}

\begin{lemma}[Sign of the Hessian]
\lean{SpinGlass.hessian\_free\_energy\_std\_basis\_offdiag\_nonpos}\leanok
For $H\in\R^{\Sigma}$ and $\gamma\ne\gamma'$, $\partial_\gamma\partial_{\gamma'}\log Z(H)=-G(\gamma)G(\gamma')\le0$,
where $G$ is the Gibbs weight.
\end{lemma}

\begin{lemma}[Derivative of $\Phi$]\label{lem:dPhi}
\lean{(to be added to ParisiFormula/GuerraToninelli.lean)}
For $t\in(0,1)$,
\[
\Phi'(t)=\frac12\sum_{\gamma,\gamma'}\big(C_L-C_{\mathrm{blk}}\big)(\gamma,\gamma')\;
\E\big[\partial_\gamma\partial_{\gamma'}\log Z(K_t)\big].
\]
\end{lemma}
\begin{proof}[Off-path proof outline; not yet formalised]
Differentiate under $\E$ by dominated convergence, apply Gaussian integration by
parts, and reduce to the covariance--Hessian trace. The files in \texttt{port/}
are historical proof-structure references, not compiled implementations of this
step; see \texttt{port/README.md} for the incompatible fork APIs. The local Stein
and cascade lemmas now prove the distinct Theorem~2.1 interpolation below, but
the Guerra--Toninelli specialization has not been assembled.
\end{proof}

\begin{proposition}[Monotonicity]\label{prop:mono}
\lean{SpinGlass.Targets.\ensuremath{\Phi}\_monotoneOn}
\uses{lem:cov, lem:dPhi}
$\Phi$ is non-decreasing on $[0,1]$.
\end{proposition}
\begin{proof}
By Lemma~\ref{lem:dPhi} and the two sign lemmas, every term of $\Phi'(t)$ is a product of
two non-positive numbers (or vanishes on the diagonal), hence $\Phi'\ge0$ on $(0,1)$;
conclude with continuity of $\Phi$ on $[0,1]$ (or the mean value theorem on compact
subintervals).
\end{proof}

\begin{theorem}[Guerra--Toninelli superadditivity]
\lean{SpinGlass.guerra\_toninelli\_superadditive}
\uses{prop:mono}
$(N+M)F_{N+M}\ge NF_N+MF_M$.
The checked theorem takes monotonicity as a hypothesis; the unconditional result
remains dependent on Proposition~\ref{prop:mono}.
\end{theorem}

\begin{theorem}[Thermodynamic limit]
\lean{SpinGlass.Targets.free\_entropy\_tendsto, SpinGlass.free\_entropy\_tendsto\_of\_bddAbove}
\uses{lem:annealed, prop:mono}
$\lim_{N\to\infty}F_N(\beta,h)$ exists.
\end{theorem}
\begin{proof}
$N\mapsto NF_N$ is superadditive and $F_N$ is bounded above (Lemma~\ref{lem:annealed});
Fekete's lemma (\texttt{Subadditive.tendsto\_lim} in Mathlib) applied to $-NF_N$ gives the
limit. This assembly step is already formalised with monotonicity and boundedness
hypotheses; \texttt{free\_entropy\_tendsto} itself remains a placeholder.
\end{proof}

\section{Milestone 2: the finite-step Parisi functional}

\begin{definition}[RSB scheme]
\lean{SpinGlass.Targets.RSBScheme}\leanok
A $k$-step scheme is a pair of sequences
$0=m_0\le m_1\le\dots\le m_k\le m_{k+1}=1$ and
$0=q_0\le q_1\le\dots\le q_{k+1}\le q_{k+2}=1$.
\end{definition}

\begin{definition}[Gaussian smoothing operator]
\lean{SpinGlass.Targets.parisiStep, SpinGlass.Parisi.T}\leanok
For $m>0$, $v\ge0$ and $A:\R\to\R$ of moderate growth,
\[
T_{m,v}A(x)=\frac1m\log\int\exp\big(mA(x+\sqrt v z)\big)\,d\gamma(z),
\qquad T_{0,v}A(x)=\int A(x+\sqrt vz)\,d\gamma(z).
\]
The zero-mass branch belongs to \texttt{parisiStep}. The ported \texttt{Parisi.T}
agrees with it only for $m\ne0$.
\end{definition}

\begin{lemma}[Semigroup law]
\lean{SpinGlass.Parisi.T\_add}\leanok
\lean{SpinGlass.Parisi.T\_add\_of\_hasLinearGrowth}\leanok
$T_{m,v_1}T_{m,v_2}A=T_{m,v_1+v_2}A$ for $m\ne0$, $v_1,v_2\ge0$, and measurable
$A$ of linear growth. The bounded-function version is compiled in
\texttt{ParisiFormula/ParisiOperator.lean}; the linear-growth extension is in
\texttt{ParisiFormula/ParisiOperatorGrowth.lean}.
\end{lemma}

\begin{definition}[Finite-step Parisi functional]
\lean{SpinGlass.Targets.parisiF, SpinGlass.Targets.parisiFunctional}\leanok
Set $F_{k+2}(x)=\log\cosh x$ and, for $p=k+1,\dots,1$,
$F_p=T_{m_p,\,\beta^2(q_{p+1}-q_p)}F_{p+1}$; finally $F_0=T_{0,\beta^2q_1}F_1$.
Then
\[
\mathcal P_k(m,q)=\log2+F_0(h)-\frac{\beta^2}{4}\sum_{p=1}^{k+1}m_p\big(q_{p+1}^2-q_p^2\big).
\]
\end{definition}

\begin{lemma}[Replica-symmetric case]
\lean{SpinGlass.Targets.parisiFunctional\_rsScheme}\leanok
For $k=0$, $m=(0,1)$, $q=(0,q,1)$:
$\mathcal P_0=\log2+\E\log\cosh(\beta\sqrt q\,z+h)+\frac{\beta^2}4(1-q)^2$.
\end{lemma}
\begin{proof}
$F_1(x)=\log\E\cosh(x+\beta\sqrt{1-q}\,z)=\log\cosh x+\frac{\beta^2(1-q)}2$, so
$F_0(h)=\E\log\cosh(\beta\sqrt qz+h)+\frac{\beta^2(1-q)}2$, and
$\frac{\beta^2(1-q)}2-\frac{\beta^2}4(1-q^2)=\frac{\beta^2}4(1-q)^2$.
This checked identity provides a sanity check on the normalisations.
\end{proof}

\begin{lemma}[Fixed-level continuity and a minimizer]
\lean{SpinGlass.Targets.continuousOn\_parisiFunctionalRaw}\leanok
\lean{SpinGlass.Targets.exists\_minimizer\_parisiFunctional}\leanok
For fixed $k$, $\beta$, and $h$, the functional is continuous in $(m,q)$ on the
compact, nonempty set of admissible schemes. It therefore attains a minimum
at that level. This is the regularity used in the current final deduction.
\end{lemma}

\begin{lemma}[Legacy uniform Lipschitz target; open and off-path]
\lean{SpinGlass.Targets.parisiFunctional\_lipschitz}
The retained target asserts a Lipschitz estimate in the $\ell^1$ distance between
scheme parameters, with a constant $C(\beta,h)$ uniform in $k$.
This stronger estimate is not a hypothesis of the current Parisi-formula deduction.
\end{lemma}

\section{Milestone 3: Guerra's upper bound}

\begin{definition}[Finite-cascade interpolation]
\lean{SpinGlass.Targets.guerraPhi, SpinGlass.Targets.guerraPsi}\leanok
Let $(z_p^i)_{0\le p\le k+1,\;i\le N}$ be independent standard Gaussians and
\[
H_t(\sigma)=\sqrt t\,U(\sigma)+\sum_{i\le N}\sigma_i
\left(h+\sqrt{1-t}\sum_{p=0}^{k+1}z^i_p\sqrt{\beta^2(q_{p+1}-q_p)}\right).
\]
Apply the finite backward log-exponential recursion to the log partition function,
using ordinary expectation at zero mass, and then average over the SK disorder
and divide by $N$. This defines $\varphi_N(t)$ in \texttt{guerraPhi}.
For $N\ge1$, the proved endpoint identities give $\varphi_N(1)=F_N$ and
$\varphi_N(0)=\log2+F_0(h)$. Set
\[
C_s=\frac{\beta^2}{4}\sum_{p=1}^{k+1}m_p(q_{p+1}^2-q_p^2),
\qquad \psi_s(t)=\log2+F_0(h)-tC_s .
\]
Thus $\psi_s(1)=\mathcal P_k(m,q)$.
\end{definition}

\begin{theorem}[Talagrand Theorem 2.1, exact-covariance SK version]\label{thm:identity}
\lean{SpinGlass.Targets.guerra\_identity}\leanok
For $N\ge1$ and every finite scheme, $\varphi_N$ is continuous on $[0,1]$ and
\[
\varphi_N'(t)=-C_s-\operatorname{Rem}_N(t), \qquad 0<t<1,
\qquad 0\le\operatorname{Rem}_N(t)\le\beta^2.
\]
The explicit witness \texttt{guerraRemainder} is $\beta^2/4$ times the
mass-weighted conditional two-replica average of the squared-overlap errors
$(R_{1,2}-q_\ell)^2$. It is extended by zero outside $(0,1)$; no endpoint
continuity of the remainder is asserted.
\end{theorem}

\begin{proof}[Implemented proof structure]
The local cascade modules prove endpoint factorization, closed-interval
continuity, and first and second tilted derivative rules. Coordinate Gaussian
integration by parts handles the disorder and site fields. The normalized
replica averages identify the covariance and field contractions; telescoping
and completion of the square give the identity. Nonnegativity, $|R-q_\ell|\le2$,
and total mass one give the remainder bounds. This proof is in
\texttt{Targets/Talagrand.lean}, not in the historical \texttt{port/} files.
\end{proof}

\begin{theorem}[Guerra's RSB bound \cite{Gue03}]\label{thm:guerra}
\lean{SpinGlass.Targets.guerra\_rsb\_bound}\leanok
\uses{thm:identity}
For every $N\ge1$, every $(\beta,h)$ and every $k$-step scheme,
$F_N(\beta,h)\le\mathcal P_k(m,q)$.
\end{theorem}

\begin{proof}
The derivative of $\varphi_N-\psi_s$ is nonpositive on $(0,1)$.
Continuity and the mean value theorem extend the comparison to the endpoints,
giving $\varphi_N(1)\le\psi_s(1)=\mathcal P_k(m,q)$.
\end{proof}

\begin{theorem}[Upper bound in the limit]
\lean{SpinGlass.Targets.limsup\_free\_entropy\_le\_parisiValue}\leanok
\uses{thm:guerra}
$\limsup_N F_N\le\mathcal P(\beta,h):=\inf_k\inf_{(m,q)}\mathcal P_k(m,q)$.
The Lean statement uses eventual $\varepsilon$ upper bounds, without assuming
existence of a limit. Nonemptiness and boundedness below of the defining set of
\texttt{parisiValue} are proved in \texttt{Targets/Milestones.lean}.
\end{theorem}

\section{Milestone 4: Theorem 2.2 is formalised}

\paragraph{Route decision (fixed).}  We formalise \textbf{Talagrand's} proof
\cite{Tal06}.  Panchenko's route \cite{Pan13} is out of scope for now.  Consequently the
following are \emph{not} prerequisites: the Ghirlanda--Guerra identities, ultrametricity of
the overlap, Ruelle probability cascades, and the Aizenman--Sims--Starr scheme.
The fixed-level minimizer used below is already proved; the legacy uniform
Lipschitz target is not a prerequisite of this deduction.

\begin{theorem}[Talagrand Theorem 2.2; formalised]\label{thm:two-two}
\lean{SpinGlass.Targets.talagrand\_theorem\_2\_2}\leanok
Fix $\beta>0$, $h\in\R$, and a family of SK disorders. For every $t_0<1$ there
exists $\varepsilon>0$ such that, whenever a scheme $s$ satisfies
\[
\mathcal P_k(s)\le\mathcal P(\beta,h)+\varepsilon,
\qquad \mathcal P_k(s)\le\mathcal P_k(s')\quad\text{for all schemes $s'$ at level $k$},
\]
one has $\varphi_N(t)\to\psi_s(t)$ for every $0\le t\le t_0$.
The theorem is checked under exactly these quantifiers and has no placeholder
dependency.
\end{theorem}

\begin{theorem}[Theorem 2.2 downstream integration; formalised]
\lean{SpinGlass.Targets.talagrand\_theorem\_2\_2\_from\_theorem\_2\_4}\leanok
Under the hypotheses and with the conclusion of Theorem~\ref{thm:two-two},
$\varphi_N(t)\to\psi_s(t)$ for every $0\le t\le t_0$. This theorem has the exact
mathematical type of the canonical target. It applies the completed Theorem~2.4
to the strict mass/overlap reduction, which already contains Proposition~2.3 and
the concentration-to-convergence deduction. The canonical public theorem now
reuses this result directly.
\end{theorem}

\begin{lemma}[Concentration-to-convergence implication]\label{lem:concentration-convergence}
\lean{SpinGlass.Targets.guerraPhi\_uniform\_of\_overlap\_concentration}\leanok
Fix a finite scheme and $0\le t_0<1$. Write
$g_N(t)=\psi_s(t)-\varphi_N(t)$, and let $T_N(t,a)$ be the mass-weighted
two-replica probability of $(R_{1,2}-q_\ell)^2\ge a$, implemented as
\texttt{guerraOverlapTail}.
Assume there is a fixed $K\ge0$ such that, for every $\eta>0$ and all
sufficiently large $N$, uniformly for $0<t<t_0$,
\[
T_N\bigl(t,Kg_N(t)+\eta\bigr)\le\eta.
\]
Then $\varphi_N\to\psi_s$ uniformly on $[0,t_0]$.
This is a proved implication. The displayed hypothesis is now supplied downstream
from Theorem~2.4 by the finite-overlap and exact-scheme-reduction modules.
\end{lemma}

\begin{proof}[Formalised deduction corresponding to (2.20)--(2.22)]
The actual gap is continuous, nonnegative, and zero at $t=0$, with
$g_N'=\operatorname{Rem}_N$ in the interior. Normalization and positivity of
the replica averages, together with $(R-q_\ell)^2\le4$, give
\[
\operatorname{Rem}_N(t)\le\frac{\beta^2K}{4}g_N(t)
  +\frac{5\beta^2}{4}\eta.
\]
The integrating-factor estimate in \texttt{guerraGap\_le\_mul\_exp} uses
only interior derivatives and endpoint continuity. Taking $\eta$ arbitrarily
small proves uniform convergence. The theorem
\texttt{talagrand\_theorem\_2\_2\_of\_overlap\_concentration} assembles
the near-optimality and fixed-level-minimality quantifiers under the explicit
concentration hypothesis. The new downstream integration theorem discharges that
hypothesis through Theorem~2.4 and has no placeholder dependencies.
\end{proof}

\begin{theorem}[Parisi formula; formalised \cite{Tal06}]
\lean{SpinGlass.Targets.parisi\_formula}\leanok
\uses{thm:identity, thm:two-two, thm:guerra}
For $\beta>0$, $\lim_N F_N(\beta,h)=\mathcal P(\beta,h)$.
The theorem is checked without a critical-path placeholder dependency.
\end{theorem}

\begin{proof}[Checked deduction from Theorem~2.2]
Theorem~2.1 gives a uniform bound $|\varphi_N'|\le L$. Choose $t_0<1$ close to
one, an $\varepsilon$-optimal scheme, and then a minimizer at the same level.
Theorem~2.2 gives $\varphi_N(t_0)\to\psi_s(t_0)\ge\mathcal P$.
The derivative bound yields
$\varphi_N(1)\ge\varphi_N(t_0)-L(1-t_0)$, hence the lower bound as $t_0\uparrow1$.
Combine this with the proved upper bound.
\end{proof}

\begin{lemma}[Talagrand Lemma 2.7: unrestricted coupled cascade]\label{lem:coupled-cascade}
\lean{SpinGlass.Targets.coupledPhi\_eq\_two\_guerraPhi}\leanok
\lean{SpinGlass.Targets.coupledObservable\_eq\_replicaMeasure}\leanok
Use independent increments at inner levels and shared increments at outer levels,
halving the mass at each shared level. For the unrestricted two-replica cascade,
\[
\frac1N\mathbb E J_1=2\varphi_N(t),\qquad
\mathbb E\bigl(V_1\cdots V_k\langle f\rangle\bigr)=\mu_r(f).
\]
The pressure identity is checked for $0\le t\le1$; the observable identity is
checked for $0<t<1$ in iterated conditional-expectation form.
Here the displayed indices follow the paper, whose $k$ is the local $k+1$.
\end{lemma}

\begin{proof}[Formalised proof]
The joint partition sum factors. Integration against the product Gaussian measure preserves
$J_\ell=F_\ell^1+F_\ell^2$ at independent levels; on the shared diagonal,
$T_{m/2}(2F)=2T_m(F)$. The corresponding normalized densities are
$V_\ell=W_\ell^1W_\ell^2$ at independent levels and $V_\ell=W_\ell$ at shared levels.
Iterating these identities proves the assertions. Zero masses are included;
finite sums are exchanged with integrals only after proving integrability.
\end{proof}

\begin{lemma}[Replica decomposition of the existing remainder]
\lean{SpinGlass.Targets.guerraReplicaExpectation\_eq\_sum\_measures}\leanok
\lean{SpinGlass.Targets.guerraOverlapTail\_le\_of\_replicaMeasure}\leanok
Use local reversed-depth indexing. Write $a_d$ for the remaining mass at depth $d$
and $\mu^{(d)}$ for the normalized replica expectation at that depth. Then
\[
\texttt{guerraReplicaExpectation}(K)
 =\sum_{d=0}^{k+1}(a_d-a_{d+1})\,\mu^{(d)}(K_d).
\]
The weights are nonnegative and sum to one; the depth-zero weight vanishes.
Thus a bound by $\delta$ on each genuine per-level overlap tail implies
$T_N(t,a)\le\delta$, the precise input to Lemma~\ref{lem:concentration-convergence}.
This connection is proved in \texttt{Targets/ReplicaMeasure.lean};
\texttt{Targets/TalagrandOverlapTail.lean} obtains the per-level bounds from
Theorem~2.4 and Proposition~2.5.
\end{lemma}

\begin{lemma}[Deterministic comparison in Talagrand Lemma 2.6]
\lean{SpinGlass.Targets.constrainedBase\_eq\_add\_log\_event}\leanok
\lean{SpinGlass.Targets.log\_coupledEvent\_le\_gap}\leanok
Assume the overlap $u$ is attainable and $m_1>0$. At each positive-mass level,
write $Q_\ell=\mathbb E_\ell(V_\ell\cdots V_k\langle 1_{R_{1,2}=u}\rangle)$
in the paper's indexing. Then
\[
0<Q_\ell\le \exp\bigl(n_\ell(J_{\ell,u}-J_\ell)\bigr)\le1,
\qquad
\log Q_\ell\le n_\ell(J_{\ell,u}-J_\ell).
\]
The implementation uses the actual constrained partition sum and its recursively
smoothed pressure.
\end{lemma}

\begin{proof}[Formalised deterministic induction]
At the Gibbs base the event probability is the ratio of the constrained and
unrestricted partition sums. Attainable overlap makes both sums positive,
giving $J_{k+1,u}=J_{k+1}+\log Q_{k+1}$.
If $B\le A$ and the next mass satisfies $0<m\le p$, then
$F\le e^{p(B-A)}\le e^{m(B-A)}$.
Integrating against the normalized density gives
$\mathbb E_m F\le e^{m(T_mB-T_mA)}$.
The two-field growth estimates justify all Gaussian integrals without
factorising the constrained function. Positivity of the tilted event is proved
separately before taking logarithms. The pressure-gap comparison excludes the
outer zero-mass average; positivity includes it.
\end{proof}

\begin{lemma}[Concentration in standard Gaussian coordinates]
\lean{SpinGlass.Targets.gaussianCoupledGap\_upper\_tail}\leanok
\lean{SpinGlass.Targets.gaussianCoupledEvent\_small}\leanok
Realise the SK disorder by its spectral coefficients and the shared outer field
by independent standard Gaussian coordinates $Z$. Let $N>0$, $t\in[0,1]$, and
$u$ be attainable. The gap $D(Z)=J_{1,u}(Z)-J_1(Z)$ satisfies
\[
\mathbb P\{D-\E D\ge a\}
\le \exp\!\left(-\frac{a^2}{32(1+|\beta|)^2N}\right),\qquad a\ge0.
\]
If $m_1>0$, $\varepsilon>0$, and $\E D\le-\varepsilon N$, then
\[
\E Q_1\le e^{-m_1\varepsilon N/4}
+\exp\!\left(-\frac{\varepsilon^2N}{128(1+|\beta|)^2}\right).
\]
These expectations are over the standard Gaussian coordinates. The constants
are uniform in $N,t,u$ and the replica split.
\end{lemma}

\begin{proof}[Formalised concentration and truncation]
The squared coefficient norm of each interpolated Hamiltonian is
$tN\beta^2/2+(1-t)N\beta^2q_1\le N\beta^2$.
The nonempty constrained log partition sum and every smoothing step preserve
sup-norm differences. Thus $D$ is $4(1+|\beta|)\sqrt N$-Lipschitz.
The proved Gaussian MGF inequality gives the displayed tail bound.
Split the expectation at $D=-\varepsilon N/2$, use
$Q_1\le\min(1,e^{n_1D})$ and $n_1\ge m_1/2$, and apply that tail estimate.
Joint measurability of the disorder-dependent cascade and integrability of
the gap and event are proved separately.
\end{proof}

\begin{lemma}[Lemma 2.6 in the abstract disorder formulation]
\lean{SpinGlass.Targets.coupledGaussian\_pair\_map}\leanok
\lean{SpinGlass.Targets.talagrand\_lemma\_2\_6}\leanok
For $N>0$, attainable $u$, $m_1>0$, and $0\le t\le1$, suppose
$\E J_{1,u}\le\E J_1-\varepsilon N$, where $\varepsilon>0$.
Then the original tilted event expectation satisfies
\[
\E\bigl(V_1\cdots V_k\langle1_{R_{1,2}=u}\rangle\bigr)
\le K e^{-N/K},\qquad
K=2+\frac1r,\quad
r=\min\left(\frac{m_1\varepsilon}{4},
\frac{\varepsilon^2}{128(1+|\beta|)^2}\right)>0.
\]
The constant is independent of $N,t,u$ and the replica split.
\end{lemma}

\begin{proof}[Formalised change of law]
The abstract disorder's independent Gaussian coordinates have the same product
law as the scaled standard coordinates, including zero spectral variances.
The joint-law identity preserves an independent outer field. At the last cascade
level the mass is zero and the increment is shared, so Fubini identifies the
standard-coordinate gap mean with $N(\Psi(t,u)-2\varphi(t))$ and the event mean
with the original tilted expectation. The constrained and unrestricted pressures
are proved integrable, justifying the subtraction of their expectations.
Apply the preceding Gaussian-coordinate estimate and use the displayed $K$.
\end{proof}

\begin{proposition}[Proposition 2.5 at interior times]
\lean{SpinGlass.Targets.talagrand\_proposition\_2\_5}\leanok
Under the same positivity and attainability conditions, for $0<t<1$,
\[
\Psi(t,u)\le2\varphi(t)-\varepsilon
\quad\Longrightarrow\quad
\mu_r(R_{1,2}=u)\le K e^{-N/K}.
\]
This follows from Lemmas~2.6 and~2.7. The interior-time domain is sufficient:
the formalised concentration-to-convergence implication only needs concentration
on $(0,t_0)$ and already concludes convergence on $[0,t_0]$.
\end{proposition}

\begin{lemma}[The Section 5 single-site construction]
\lean{SpinGlass.Targets.coupledSite\_spin\_sum}\leanok
\lean{SpinGlass.Targets.lambdaCoupledBase\_time\_zero}\leanok
Define
\[
A^*_{\lambda}(x,y)=\log\bigl(\cosh x\cosh y\cosh\lambda
  +\sinh x\sinh y\sinh\lambda\bigr).
\]
Its logarithm argument is strictly positive, and
\[
\sum_{a,b=\pm1}e^{ax+by+\lambda ab}=4e^{A^*_{\lambda}(x,y)}.
\]
The unrestricted interacting partition function at zero interpolation time
therefore has logarithm
\[
2N\log 2+\sum_i A^*_{\lambda}(x_i+h,y_i+h).
\]
Also $A^*_0(x,y)=\log\cosh x+\log\cosh y$, and its derivative in $\lambda$
at zero is $\tanh x\tanh y$.
\end{lemma}

\begin{proposition}[Lambda comparison for the existing coupled cascade]
\lean{SpinGlass.Targets.constrainedCascade\_le\_lambda}\leanok
\lean{SpinGlass.Targets.constrainedPhi\_le\_lambdaCoupledPhi}\leanok
Let $P_{\lambda}(t)$ denote the normalized expected pressure of the existing
coupled cascade with terminal energy
$H_t(\sigma)+H_t(\tau)+\lambda\sum_i\sigma_i\tau_i$.
For $N>0$, attainable $u$, $0\le t\le1$, and $d\le k+1$,
\[
\Psi(t,u)\le P_{\lambda}(t)-\lambda u,\qquad
P_0(t)=2\varphi(t),\qquad
|P_{\lambda}(t)-P_{\mu}(t)|\le|\lambda-\mu|.
\]
\end{proposition}
\begin{proof}
On $R_{1,2}=u$ the extra interaction equals $\lambda Nu$, so the restricted
partition sum is at most $e^{-\lambda Nu}$ times the unrestricted interacting
sum. Monotonicity and additive-constant equivariance propagate this through
every cascade level, including zero mass. Gaussian-disorder integrability
follows from the bound $N|\lambda|$ relative to the unrestricted cascade.
Integration gives the pressure comparison; Lemma~2.7 gives the normalization.
The terminal energies change by at most $N|\lambda-\mu|$, and cascade
nonexpansiveness gives the final bound after division by $N$.
\end{proof}

These are supporting results for the endpoint argument in (5.17).
The specific left-interval second-interpolation endpoints are now proved below;
the derivative bound of Theorem~3.1 remains open. In particular $P_0=2\varphi$
does not supply Theorem~2.4's strict improvement relative to $2\psi$.
The terminal derivative above is not the full Lemma~5.8 cascade identity.

\begin{lemma}[General finite paired recursion via existing RSAT results]
\lean{SpinGlass.Targets.pairedScalarCascade\_good}\leanok
\lean{SpinGlass.Targets.pairedVectorCascade\_eq\_sum}\leanok
Fix any sequence of masses $m_j\ge0$ and continuous signed coefficients
$a_j(p),b_j(p)$, where $p$ belongs to a first-countable topological space.
Starting with $F_0(p,\lambda;x,y)=A^*_{\lambda}(x,y)$, let
\[
F_{j+1}(p,\lambda;x,y)
=\begin{cases}
\E_z F_j(p,\lambda;x+a_j(p)z,y+b_j(p)z),&m_j=0,\\
\dfrac1{m_j}\log\E_z e^{m_jF_j(p,\lambda;x+a_j(p)z,y+b_j(p)z)},&m_j>0.
\end{cases}
\]
Let $\mathcal F_j$ be the corresponding $N$-site recursion, starting with
$\sum_i A^*_{\lambda}(x_i,y_i)$ and using independent standard $z_i$ at each
step. At every finite depth,
\[
\mathcal F_j(p,\lambda;x,y)=\sum_i F_j(p,\lambda;x_i,y_i).
\]
Moreover $F_j$ and its lambda derivative are jointly continuous, the derivative
is bounded in absolute value by one, and $F_j$ is one-Lipschitz in each field.
\end{lemma}
\begin{proof}
The local terminal is proved equal to RSAT's \texttt{gtTerminal}. Iterate
RSAT's \texttt{step0\_good} and \texttt{stepM\_good}, then apply its
\texttt{gtVectorStep\_sum} at every step. The regularity property supplies
the required integrability and positivity. The old shared Gaussian step is
the general vector step at mass $m/2$; its independent step is two vector
steps at mass $m$, with coefficients $(\sqrt v,0)$ and $(0,\sqrt v)$.
\end{proof}
The lambda derivative is proved at fixed masses; neither the
mass-variation estimates nor the identification with $U'$ is asserted here.

\begin{proposition}[Second-interpolation endpoints, left-interval construction]
\lean{SpinGlass.Targets.section5Interpolation\_one}\leanok
\lean{SpinGlass.Targets.section5Interpolation\_zero\_le}\leanok
Use the paper's indices in this proposition. Let $1\le r\le k$,
$q_{r-1}\le u\le q_r$, $0\le t\le1$, and
$m_{r-1}/2\le m\le m_r$. Insert $u$ into the overlap sequence at index $r$.
Give levels below $r$ masses $m_p/2$, the inserted level mass $m$, and levels
above $r$ masses $m_{p-1}$, as in (5.5)--(5.6). Denote these sequences by
$(\rho_p)$ and $(n_p)$. Put
\[
S_p=\begin{cases}
\beta^2(q_{p+1}-q_p),&p<r,\\
0,&p=r,\\
\beta^2(q_p-q_{p-1}),&p>r,
\end{cases}
\qquad
D_p(w)=(1-t)S_p+(1-w)t\beta^2(\rho_{p+1}-\rho_p).
\]
Define $\eta(w)$ by the constrained terminal with disorder
$\sqrt{wt}\,H$, external field $h$, and nested Gaussian increments of variance
$D_p(w)$: shared below $r$ and independent at or above $r$.
The scalar $V(\lambda,m,v)$ is the finite recursion (5.10)--(5.15).
For $N>0$ and attainable overlap $u$,
\[
\eta(1)=\Psi(t,u),\qquad
\eta(0)\le2\log2+V(\lambda,m,v(u))-\lambda u,
\quad v(u)=t\beta^2(q_r-u).
\]
The inserted mass and overlap sequences are nondecreasing with the prescribed
endpoints, and all $D_p(w)$ are nonnegative on $0\le w\le1$.
\end{proposition}
\begin{proof}
At $w=1$ the inserted level has variance zero, so its logarithmic-Laplace
transform is the identity, including at mass zero. Delete this level and move
the factor $\sqrt{1-t}$ from the Gaussian increments back to the terminal
fields. An exact affine-coordinate identity through the nested integrals
recovers the original constrained cascade, proving (5.8).

At $w=0$ the disorder vanishes. The variances $D_p(0)$ are exactly those in
(5.10)--(5.13) with $v=v(u)$. The overlap penalty bounds the restricted
terminal by $2N\log2-\lambda Nu+\sum_i A^*_{\lambda}(x_i,y_i)$.
Monotonicity and constant equivariance propagate the inequality. Applying the
proved site tensorization at every level yields $N V(\lambda,m,v(u))$;
division by $N$ gives (5.17). RSAT's Gaussian regularity theorems also give
the actual lambda derivative of this $V$ at fixed $m$.
\end{proof}
These endpoint results alone do not assert the intervening covariance derivative
bound. The right-interval construction and both actual interior-time pressure
derivatives and the negative initial-overlap case are proved in the subsequent
modules. Proposition~5.7's actual tagged mixed-pressure inequality is now
checked below, together with strict scalar sorting and endpoint transport;
its remaining boundary assembly is still open. The uniform bound of
Lemma~5.9 is also checked.

\begin{lemma}[Constrained terminal derivatives and the SK covariance trace]
\lean{SpinGlass.Targets.hasDerivAt\_constrainedPairFieldBase\_second}\leanok
\lean{SpinGlass.Targets.constrainedPairSecond\_SK\_trace}\leanok
Let $S_u$ be the nonempty finite set of pairs with overlap $u$, and let $F(U)$
be the constrained log partition at fixed physical fields. If
$V^{\rm pair}(\sigma,\tau)=V(\sigma)+V(\tau)$, then
\[
D F(U)[V]=\langle V^{\rm pair}\rangle_U,\qquad
D^2 F(U)[V,W]=\operatorname{Cov}_U(V^{\rm pair},W^{\rm pair}).
\]
Write $R_{ab}(p,q)$ for the four overlaps between pair states and put
$K(p,q)=\frac{\beta^2}{2}\sum_{a,b}R_{ab}(p,q)^2$.
For the SK spectral directions $w_i$ and variances $\tau_i$, and $N>0$,
\[
\sum_i\tau_i D^2F(U)[w_i,w_i]
=N\left(\beta^2(1+u^2)-\langle K(p,q)\rangle_U^{\otimes2}\right).
\]
\end{lemma}
\begin{proof}
The constrained terminal is identified with RSAT's finite-state log partition
on $S_u$, preserving Talagrand's positive-Hamiltonian convention. Its existing
first and Gibbs-weight derivative formulas give the mixed Hessian. The local
SK spectral covariance identity supplies the four overlap squares. The diagonal
matrix of a constrained pair is $\left(\begin{smallmatrix}1&u\\u&1\end{smallmatrix}\right)$,
so its normalized covariance is the constant $\beta^2(1+u^2)$.
\end{proof}

\begin{lemma}[Covariance algebra after the intended nested IBP]
\lean{SpinGlass.Targets.pairCovarianceExpression\_eq}\leanok
\lean{SpinGlass.Targets.pairCovarianceExpression\_le}\leanok
Let $0=m_0\le\cdots\le m_\kappa=1$, $\rho_0=c_0=0$,
$\rho_{\kappa+1}=1$, and $c_{\kappa+1}=u$. Put
\[
Q_\ell=\begin{pmatrix}\rho_\ell&c_\ell\\c_\ell&\rho_\ell\end{pmatrix},\quad
L_\ell(p,q)=\beta^2\sum_{a,b}(Q_\ell)_{ab}R_{ab}(p,q),\quad
D_\ell(p,q)=\frac{\beta^2}{2}\sum_{a,b}(R_{ab}(p,q)-(Q_\ell)_{ab})^2.
\]
For arbitrary nonnegative normalized weights $\mu_\ell$ on $S_u^2$, define
\[
\mathcal B=\frac t2\left[-\beta^2(1+u^2)
 +\sum_{\ell=1}^{\kappa}(m_{\ell-1}-m_\ell)\mu_\ell(K-L_\ell)\right],
\]
\[
\mathcal C=\sum_{p=0}^{\kappa}m_p(E_{p+1}-E_p),\quad
E_p=\frac{\beta^2}{2}(\rho_p^2+c_p^2).
\]
For $t\ge0$,
\[
\mathcal B=-t\mathcal C-\frac t2\sum_{\ell=1}^{\kappa}
 (m_\ell-m_{\ell-1})\mu_\ell(D_\ell)\le-t\mathcal C.
\]
\end{lemma}
\begin{proof}
RSAT's scalar quadratic remainder identity gives
$K-L_\ell=D_\ell-2E_\ell$ entry by entry. Normalize the constant terms under
$\mu_\ell$ and apply Mathlib's summation-by-parts identity. The remaining
weighted squares are nonnegative. For $c_\ell=\eta\rho_{\min(\ell,\tau)}$,
$\eta^2=1$, the correction doubles exactly at the levels $p<\tau$.
\end{proof}

\begin{lemma}[Inserted correction in (5.9)]
\lean{SpinGlass.Targets.section5Correction\_eq}\leanok
For the inserted left-interval sequences (5.5)--(5.6), using paper indices,
\[
\mathcal C=\sum_{p=0}^{k}m_p\bigl(\theta(q_{p+1})-\theta(q_p)\bigr)
 +(m-m_{r-1})\bigl(\theta(q_r)-\theta(u)\bigr),\qquad
\theta(q)=\frac{\beta^2q^2}{2}.
\]
\end{lemma}
\begin{proof}
Below the split, the doubled shared contribution cancels the half mass.
The two split increments combine into the original increment plus the
displayed correction. Above the split, shift the index by one. The original
sum is exactly twice the existing Parisi correction.
\end{proof}
The quantity $\mathcal B$ above is now identified with the actual derivative
for both positive-overlap Section~5 constructions below. The covariance
inequality and endpoint transport give (5.9) and its dual. The general
signed-field version of Theorem~3.1 and Theorem~2.4 remain open.

\begin{lemma}[Fixed-variance differentiation through the paired cascade]
\lean{SpinGlass.Targets.hasDerivAt\_constrainedPairFieldCascade}\leanok
\lean{SpinGlass.Targets.constrainedPairFieldCascadeD\_abs\_le}\leanok
For fixed nonnegative masses and variances, the actual constrained cascade
with disorder $U+aV$ is differentiable in $a$ at every finite depth.
Its derivative is the iterated normalized tilt of the terminal Gibbs mean
of $V(\sigma)+V(\tau)$ and has absolute value at most
$2\max_\sigma|V(\sigma)|$, without a depth factor.
At one second-field or shared step, the tilted observable derivative is
\[
 \partial_a\mathbb E_W G
 =\mathbb E_W G'+m\bigl(\mathbb E_W(GA')-\mathbb E_WG\,\mathbb E_WA'\bigr).
\]
The actual one-shared-level mixed Hessian is proved as a specialization.
\end{lemma}
\begin{proof}
Reuse the existing one-field parameter chain rules. An independent level
is two transforms by Fubini; a shared level is one diagonal-shift transform
at its actual mass. Propagate measurable growth and the derivative bound
together, starting from the already proved finite-state calculus.
\end{proof}
The following extension supplies the full fixed-variance disorder Hessian and Stein
calculation, but does not yet differentiate the varying-variance interpolation.

\begin{lemma}[Full-depth mixed disorder Hessian and Gaussian Stein]
\lean{SpinGlass.Targets.hasDerivAt\_constrainedPairFieldCascade\_second}\leanok
\lean{SpinGlass.Targets.stein\_constrainedPairFieldCascadeTrace}\leanok
Let $F_j(U)$ be the actual constrained paired cascade at depth $j$, with
fixed nonnegative masses and variances. Its actual mixed disorder derivative obeys
\[
 |D^2_{V,W}F_j(U)|\le
 \left(8+16\sum_{\ell<j}|m_\ell|\right)
 \max_\sigma|V(\sigma)|\max_\sigma|W(\sigma)|.
\]
For Gaussian disorder $Z=\sum_i c_i w_i$ with spectral variances $\tau_i$,
\[
 \mathbb E\sum_i c_i D_{w_i}F_j(Z)
 =\mathbb E\sum_i\tau_i D^2_{w_i,w_i}F_j(Z).
\]
All integrability and disorder measurability are proved.
\end{lemma}
\begin{proof}
Iterate the existing tilted-observable covariance rule: twice per independent
level and once per shared level. The uniform Hessian bound implies continuity
of first directions by Mathlib's mean value theorem. Measurable line derivatives
give disorder measurability of the Hessian. Apply the existing coordinate Stein
theorem and exchange the finite spectral sum with expectation.
\end{proof}

\begin{lemma}[Separate spatial derivatives through the full paired cascade]
\lean{SpinGlass.Targets.hasDerivAt\_constrainedPairFieldCascade\_field\_second}\leanok
For a field direction $(A,B)$, the terminal observable is
$\sum_i A_i\sigma_i+\sum_i B_i\tau_i$. Its actual first derivative and the
mixed derivative with another direction $(C,D)$ propagate through every level.
Their absolute values are at most $\|A\|_1+\|B\|_1$ and
\[
 \left(2+4\sum_{\ell<j}|m_\ell|\right)
 (\|A\|_1+\|B\|_1)(\|C\|_1+\|D\|_1),
\]
respectively. Both are measurable and uniform in the fields and disorder.
\end{lemma}
\begin{proof}
Embed the separate directions as finite-state spin observables and reuse RSAT's
terminal differential calculus. Field translations commute with every Gaussian
step, identifying the propagated parameter derivatives with spatial derivatives.
Apply the already proved full-depth first and mixed derivative inductions.
\end{proof}

\begin{lemma}[Scalar mass variation and the zero-lambda baseline]
\lean{SpinGlass.Targets.hasDerivAt\_parisiStep\_mass\_entropy}\leanok
\lean{SpinGlass.Targets.section5V\_zero\_eq\_two\_section4T}\leanok
\lean{SpinGlass.Targets.section5V\_zero\_baseline}\leanok
For a measurable linear-growth input $A$ held fixed, write
$B_m(x)=m^{-1}\log\mathbb E e^{mA(x+\sqrt v g)}$ when $m\ne0$
and let $W_m$ be its normalized exponential density. Then
\[
 \partial_m B_m(x)=\frac{\mathbb E[W_m A]-B_m(x)}m
 =\frac{\mathbb E[W_m\log W_m]}{m^2}\ge0.
\]
For the actual split scalar recursion $T(v,m)$ of (4.35),
\[
 V(0,m,v)=2T(v,m),\qquad T(v,m_{r-1})=A_0(h),\qquad
 V(0,m_{r-1},v)=2A_0(h).
\]
The baseline identities include zero masses and both variance endpoints.
\end{lemma}
\begin{proof}
Mathlib's moment-generating-function calculus applies because the existing
Gaussian growth estimates give all exponential moments. Normalization and
$w\log w\ge w-1$ give the entropy sign. For the zero-lambda identities,
reuse independent-step factorization and diagonal mass cancellation.
The scalar Gaussian semigroup merges the two split increments at equal
mass; its zero-mass expectation branch follows from Gaussian convolution.
\end{proof}
The nested mass estimates, Proposition~4.6 and stationarity after exact
scheme reduction are proved below. Uniform higher variance regularity is open.
The identification with $U'$ is checked by the following bridge.

\begin{lemma}[The actual Lemma 5.8 identity]
\lean{SpinGlass.Targets.deriv\_section5V\_zero\_eq\_Q}\leanok
\lean{SpinGlass.Targets.talagrand\_lemma\_5\_8}\leanok
\lean{SpinGlass.Targets.talagrand\_lemma\_5\_8\_within}\leanok
\lean{SpinGlass.Targets.section5Interpolation\_zero\_lambda\_gain\_Q}\leanok
For $m\ge0$ and $0\le v\le a$,
\[
 \left.\partial_\lambda V(\lambda,m,v)\right|_{\lambda=0}=Q(m,v).
\]
At baseline mass $m_{r-1}<1$ this equals the ordinary $U'(v)$ for
$0<v<a$, and its inward derivative at the endpoints. Numerical uniqueness
of the latter is asserted for $a>0$; the derivative predicate remains valid
for a degenerate interval. The optimized time-zero gain therefore has the
explicit deficit $(Q(m_{r-1},v(u))-u)^2/2$.
\end{lemma}
\begin{proof}
The actual lambda-derivative recursion propagates the product of scalar
slopes through independent levels at lambda zero. On the diagonal the
shared levels see twice the scalar potential and half the scalar mass,
so their normalized means agree exactly with the recursion defining $Q$.
Use the proved $U'=Q$ and the inward FTC at baseline. Substitute into
the already checked lambda-curvature gain; the $Q$ gain itself also allows
baseline mass one without claiming the $U'$ identity there.
\end{proof}
The gain is transported to the actual constrained free energy below; its
uniform comparison with $(u-q_r)^2$ near $q_r$ remains open.

\begin{lemma}[The scalar heat equation (4.4) on actual Parisi inputs]
\lean{SpinGlass.Targets.hasDerivAt\_parisiStep\_parisiF\_variance}\leanok
\lean{SpinGlass.Targets.hasFDerivAt\_parisiStep\_variance\_spatial}\leanok
If $A$ is any actual Parisi recursion input and
$B(v,x)=\operatorname{ParisiStep}(m,v,A)(x)$, then, for $v>0$,
\[
 \partial_v B=\tfrac12(B_{xx}+mB_x^2).
\]
The formula includes $m=0$. For $m\in[0,1]$, the actual first and second
spatial derivatives satisfy $0\le B_{xx}\le1-B_x^2$, and hence
$0\le\partial_v B\le\tfrac12$. Joint variance/spatial differentiation and
the chain rule for $B(v,x+\sqrt{a-v}\,z)$ hold on $0<v<a$.
\end{lemma}
\begin{proof}
Apply Gaussian Stein to the amplitude derivative and use exponential-growth
domination. The existing spatial smoothing invariant propagates through all
Parisi levels, including mass zero. RSAT's general Gaussian transform gives
joint continuity of $B$, $B_x$, $B_{xx}$; Mathlib's continuous-partials theorem
then gives the joint derivative needed by the moving-field chain rule.
\end{proof}

\begin{lemma}[The two-step split-variance identity (4.11)]
\lean{SpinGlass.Targets.hasDerivAt\_split\_parisiF}\leanok
Let $A$ be any actual Parisi input and put
\[
 \begin{aligned}
 B(v,x)&=\operatorname{ParisiStep}(m,v,A)(x),\\
 C(v,x)&=\operatorname{ParisiStep}(m',a-v,B(v,\cdot))(x).
 \end{aligned}
\]
For $m\in[0,1]$, $m'\in\mathbb R$, and $0<v<a$,
\[
 \partial_v C(v,x)=\frac{m-m'}2\,\mathbb E_W
       [B_x(v,x+\sqrt{a-v}\,g)^2],
\]
where $W$ is the actual outer normalized tilt. Both zero-mass cases are included.
The derivative is nonnegative if $m'\le m$.
\end{lemma}
\begin{proof}
Adapt the existing local-neighborhood Gaussian parameter chain rule to one
coordinate. The joint variance/spatial derivative supplies the moving-inner-field
derivative. Normalized scalar Stein cancels the two $B_{xx}$ contributions,
leaving exactly the displayed squared-first-derivative term.
\end{proof}
\begin{lemma}[Closed-interval comparison and full scalar variance derivative]
\lean{SpinGlass.Targets.monotoneOn\_section4T}\leanok
\lean{SpinGlass.Targets.hasDerivAt\_section4T\_variance}\leanok
\lean{SpinGlass.Targets.abs\_deriv\_section4T\_variance\_le}\leanok
For $m_{r-1}\le m\le1$, the actual $T(v,m)$ is continuous and nondecreasing
on $[0,a]$, with $T(0,m)=A_0(h)$. On $0<v<a$, its variance derivative
is the nested normalized mean of the derivative in (4.11), and
\[
 |\partial_v T(v,m)|\le (m-m_{r-1})/2.
\]
The bound is independent of the number of outer levels, including those of mass zero.
On the closed interval this also gives
$|T(v,m)-T(w,m)|\le(m-m_{r-1})|v-w|/2$.
\end{lemma}
\begin{proof}
RSAT supplies joint continuity of the split transform. Combine the interior
derivative sign with Mathlib's closed-interval monotonicity theorem. The existing
Gaussian order-preservation and parameter-derivative rules propagate the
comparison, derivative and bound through every unchanged outer scalar level.
\end{proof}
Uniform higher variance regularity is open; stationarity at every level
after exact scheme reduction is proved below.

\begin{lemma}[Analytic extension through zero mass]
\lean{SpinGlass.Targets.analyticAt\_parisiStep\_mass}\leanok
\lean{SpinGlass.Targets.hasDerivAt\_parisiStep\_mass\_zero}\leanok
For measurable $A$ of linear growth, put $X=A(x+\sqrt v\,g)$ with $g$ standard
Gaussian. The actual Parisi step, including its expectation branch, is analytic
in every real mass, and
\[
 \left.\partial_m\operatorname{ParisiStep}(m,v,A)(x)\right|_{m=0}
   =\tfrac12\mathbb E[(X-\mathbb E X)^2].
\]
Together with the entropy sign, this proves mass monotonicity on all of $\mathbb R$.
\end{lemma}
\begin{proof}
All exponential moments exist by the Gaussian linear-growth estimates. Identify
the actual step with the divided difference at zero of the cumulant-generating
function. Reuse Mathlib's analytic divided-difference and second CGF derivative
theorems; the expectation branch agrees exactly with the derivative at zero.
\end{proof}

\begin{lemma}[Actual nested baseline mass derivative and first variation]
\lean{SpinGlass.Targets.hasDerivAt\_section4T\_mass\_baseline}\leanok
\lean{SpinGlass.Targets.section4U\_zero\_baseline}\leanok
\lean{SpinGlass.Targets.hasDerivAt\_section4Phi\_mass\_baseline}\leanok
For $1\le r\le k+1$ and $0\le v\le a=\beta^2(q_r-q_{r-1})$, the actual
full scalar recursion is differentiable in its inserted mass at $m_{r-1}$.
Thus the function of (4.42) is genuinely defined by
\[
 U(v)=2\left.\partial_m T(v,m)\right|_{m=m_{r-1}},\qquad U(0)=0.
\]
At a zero baseline, let $B(m,v,x)$ be the actual inserted scalar step.
Then every unchanged outer mass is zero and
\[
 T(v,m)=\mathbb E_g B(m,v,h+\sqrt{\beta^2q_r-v}\,g).
\]
In this case $U(v)$ is the outer average of the centered second moment of
the actual inserted input. For the actual functional $\Phi$ of (4.37),
\[
 \left.\partial_m\Phi(m,u)\right|_{m=m_{r-1}}
 =f(u):=\tfrac12U\bigl(\beta^2(q_r-u)\bigr)
        -\frac{\beta^2}4(q_r^2-u^2),\qquad f(q_r)=0.
\]
Both overlap endpoints, coincident overlaps and zero baseline mass are included.
\end{lemma}
\begin{proof}
For positive mass, the existing tilted first-moment bound gives domination
uniform in the field on a compact positive mass interval. Propagate the
actual derivative through the normalized outer parameter rules. At zero mass,
Herbst and reflection of input and mass give the two-sided anchored bound
$|B(m,v,x)-B(0,v,x)|\le |m|v/2$. Scheme monotonicity forces every outer mass
to vanish, and the Gaussian semigroup collapses their variance to
$\beta^2q_r-v$. Mathlib's anchored dominated-differentiation theorem then
passes the actual scalar mass-zero derivative through this expectation.
Zero variance is constant in mass. Finally differentiate the affine correction
in (4.37), obtaining exactly (4.46).
\end{proof}
This proves first mass differentiation. The identities for $U'$ and $U''$
and uniform first/second mass bounds, together with Proposition~4.6, are
proved below. Stationarity is completed after exact scheme reduction;
Proposition~4.10 is completed below using proved uniform Hessian-square regularity.

\begin{lemma}[Uniform scalar mass bounds and the first half of Lemma 4.5]
\lean{SpinGlass.Targets.parisiStep\_parisiF\_mass\_derivatives\_uniform}\leanok
\lean{SpinGlass.Targets.section4T\_mass\_derivative\_uniform}\leanok
\lean{SpinGlass.Targets.section4Phi\_mass\_derivative\_uniform}\leanok
For every actual Parisi input $A$, $0\le m\le1$ and $0\le v\le V$,
\[
 |\partial_m T_{m,v}A(x)|\le K_1(V),\qquad
 |\partial_m^2 T_{m,v}A(x)|\le K_2(V),
\]
with explicit finite constants independent of $x$ and the input depth.
For the actual full split recursion and inserted functional, on their closed
physical variance/overlap intervals and $m_{r-1}\le m\le1$,
\[
 |\partial_m T(v,m)|\le K_1(\beta^2),\qquad
 |\partial_m\Phi(m,u)|\le K_1(\beta^2)+\beta^2/4.
\]
The statements include actual differentiability at mass zero.
\end{lemma}
\begin{proof}
Center the Gaussian input at $A(x)$. Its first three tilted moments have
field-independent bounds because $A$ is 1-Lipschitz. Mathlib's analytic
cumulant derivatives and the actual zero-mass extension give, with
$M(m)=T_{m,v}A(x)$ and $K(m)=mM(m)$,
\[
 (m^2M')'=mK'',\qquad (m^3M'')'=m^2K'''.
\]
Integrating these weighted identities removes the apparent singularities
at zero mass. Normalized outer means preserve the first bound without a
depth factor; the already checked full mass derivative and affine correction
give the two nested first bounds. The next lemma completes the nested second
derivative bound in Lemma~4.5.
\end{proof}

\begin{lemma}[Full nested second mass bound and baseline Taylor estimate]
\lean{SpinGlass.Targets.section4T\_second\_mass\_derivative\_uniform}\leanok
\lean{SpinGlass.Targets.section4Phi\_second\_mass\_derivative\_uniform}\leanok
\lean{SpinGlass.Targets.section4Phi\_mass\_taylor\_baseline}\leanok
On the full physical variance/overlap interval, for every $m\in[0,1]$,
the genuine second mass derivatives of $T$ and $\Phi$ are bounded in absolute
value by $C(\beta)=K_2(\beta^2)+K_1(\beta^2)^2$. In particular, writing
$f(u)=\partial_m\Phi(m_{r-1},u)$, for $m_{r-1}\le m\le1$,
\[
 |\Phi(m,u)-\mathcal P_k(m_\bullet,q_\bullet)-f(u)(m-m_{r-1})|
 \le C(\beta)(m-m_{r-1})^2.
\]
The constant is independent of depth, field and adjacent mass gaps.
\end{lemma}
\begin{proof}
Reflection and dilation of the already controlled scalar input give uniform
local mass domination on a two-sided neighborhood of $[0,1]$. Apply the
existing normalized observable derivative rule twice. If $D,E$ denote the
inner first and second mass derivatives, the invariant
$|E|+D^2\le K_2+K_1^2$ is preserved by every fixed outer mass in $[0,1]$;
the variance term is absorbed by the squared first derivative. Thus no factor
accumulates with depth. The correction defining $\Phi$ is affine in mass.
Two applications of the mean-value theorem give the displayed Taylor bound,
without optimizing its constant to $C/2$. The checked baseline value and
first-variation identities give the expansion. The next result applies
scheme optimality; the full improvement in Theorem~2.4 remains open.
\end{proof}

\begin{lemma}[Quantitative optimality: Proposition 4.6]
\lean{SpinGlass.Targets.section4FirstVariation\_lower\_bound}\leanok
\lean{SpinGlass.Targets.exists\_section4FirstVariation\_lower\_bound}\leanok
Let the scheme minimize its functional at its fixed level, and suppose
$\mathcal P_k\le\inf\mathcal P+\varepsilon$. If $m_{r-1}<m_r$, then
\[
 f(u)\ge-L(\beta)\sqrt\varepsilon\qquad(q_{r-1}\le u\le q_r),
 \quad L(\beta)=2\sqrt{C(\beta)+1}>0.
\]
The constant is independent of the level count, external field, overlap and
mass gap. Both overlap endpoints and baseline mass zero are included.
\end{lemma}
\begin{proof}
Write $C_+=C(\beta)+1$. If $f(u)<0$, the upper-mass competitor and the
Taylor bound give $-f(u)\le C_+(m_r-m_{r-1})$. Thus
$\delta=-f(u)/(2C_+)$ is an admissible mass increase. Apply near-global
minimality to this actual inserted competitor and the same Taylor bound:
$-\varepsilon\le f(u)\delta+C_+\delta^2=-f(u)^2/(4C_+)$.
Taking square roots proves the claim; $f(u)\ge0$ is immediate.
\end{proof}

\begin{lemma}[Closed-interval bounds for the actual first variation]
\lean{SpinGlass.Targets.section4U\_sub\_mem\_Icc}\leanok
\lean{SpinGlass.Targets.section4FirstVariation\_dist\_le}\leanok
If $m_{r-1}<1$, then on the full closed split interval
\[
 0\le U(w)-U(v)\le w-v\quad(v\le w),\qquad 0\le U(v)\le v.
\]
Thus $U$ is nondecreasing and 1-Lipschitz. On $[q_{r-1},q_r]$,
$|f(u)-f(w)|\le \beta^2|u-w|/2$. Zero baseline mass and coincident overlaps
are included; the strict-mass reduction supplies $m_{r-1}<1$ on the critical path.
\end{lemma}
\begin{proof}
Apply the already proved full $T$ variance comparison at mass $m_{r-1}+\delta$,
divide by $\delta>0$, and let $\delta\downarrow0$ using the actual baseline
mass derivative. The constant baseline $T(v,m_{r-1})=A_0(h)$ cancels.
Closedness of the comparison interval passes the bounds to the limit.
The formula for $f$ and $0\le u,w\le1$ give its overlap-Lipschitz bound.
No derivative of $U$ is used.
\end{proof}

\begin{lemma}[Normalized squared-slope factor and its integral identity]
\lean{SpinGlass.Targets.section4VarianceD\_eq\_massGap\_mul}\leanok
\lean{SpinGlass.Targets.section4TVarianceQ\_mem\_Icc}\leanok
\lean{SpinGlass.Targets.section4T\_sub\_eq\_massGap\_mul\_integral}\leanok
Let $Q(m,v)$ be the actual nested normalized mean of the squared spatial
slope at the split. This definition does not divide by $m-m_{r-1}$ and gives
$0\le Q(m,v)\le1$ for $0\le m\le1$, including baseline masses zero and one.
For $m_{r-1}\le m\le1$ and $0\le v\le w\le a$,
\[
 T(w,m)-T(v,m)=\frac{m-m_{r-1}}2\int_v^w Q(m,z)\,dz.
\]
The corresponding derivative identity holds at interior variance.
\end{lemma}
\begin{proof}
Linearity of each normalized mean factors the already checked full variance
derivative, and its squared-slope invariant gives the interval bound.
Mathlib's integrability and fundamental theorem for a nonnegative derivative
give the displayed identity using endpoint continuity. At equal masses both
sides vanish; no integrability of $Q$ at that baseline is inferred from the
vanishing product. For a positive mass gap, division also proves $Q$ integrable.
\end{proof}
\begin{lemma}[Baseline continuity and the actual derivative of $U$]
\lean{SpinGlass.Targets.continuousOn\_parisiStep\_parisiF\_joint}\leanok
\lean{SpinGlass.Targets.continuousOn\_section4TVarianceQ}\leanok
\lean{SpinGlass.Targets.section4U\_eq\_integral}\leanok
\lean{SpinGlass.Targets.hasDerivAt\_section4U}\leanok
\lean{SpinGlass.Targets.hasDerivAt\_section4FirstVariation\_overlap}\leanok
The actual $Q$ is continuous on $[0,1]\times[0,a]$, including both mass
endpoints and degenerate variance intervals. If $m_{r-1}<1$, then
\[
 U(v)=\int_0^v Q(m_{r-1},z)\,dz,\qquad
 U'(v)=Q(m_{r-1},v)\in[0,1]\quad(0<v<a).
\]
For $q_{r-1}<u<q_r$, the actual first variation satisfies
\[
 f'(u)=\frac{\beta^2}{2}
 \left(u-Q\bigl(m_{r-1},\beta^2(q_r-u)\bigr)\right).
\]
\end{lemma}
\begin{proof}
Dominated Gaussian continuity gives joint continuity of the scalar normalized
slope quotient; the potential's zero-mass branch follows from its Herbst gap.
Propagate the bounded squared-slope observable through the normalized means.
Continuity and $0\le Q\le1$ justify passing $m=m_{r-1}+\delta$ to baseline
in the variance integral as $\delta\downarrow0$. The actual mass derivative
identifies the left-hand quotient with $(U(w)-U(v))/2$; the right side tends
to half the baseline integral. Use $U(0)=0$ and the interior FTC. The chain
rule gives $f'$, with $\beta=0$ treated as a constant function.
\end{proof}
\begin{lemma}[Inward derivatives at the physical endpoints]
\lean{SpinGlass.Targets.hasDerivWithinAt\_section4U}\leanok
\lean{SpinGlass.Targets.derivWithin\_section4U\_eq}\leanok
\lean{SpinGlass.Targets.hasDerivWithinAt\_section4FirstVariation\_overlap}\leanok
The formulas $U'=Q$ and $f'=\beta^2(u-Q)/2$ hold as derivatives within the
closed physical intervals. In particular the occurrence of $U'(0)$ means
the inward derivative. Uniqueness of the numerical derivative within the
variance interval is asserted only when its length is positive.
\end{lemma}
\begin{proof}
Extend the continuous $Q$ to the real line by projection onto $[0,a]$ and
apply the ordinary FTC to its integral. On $[0,a]$ this integral equals the
actual $U$, so transfer the derivative back within that interval. The
variance map carries the closed overlap interval into $[0,a]$; apply the
within-set chain rule to the actual $f$. The extension is only a proof
device, not a claim about the two-sided derivative of the original $U$.
\end{proof}

\begin{lemma}[Lower-endpoint optimality for noninitial intervals]
\lean{SpinGlass.Targets.section4Phi\_lower\_overlap\_min}\leanok
\lean{SpinGlass.Targets.section4FirstVariation\_lower\_nonneg}\leanok
For $r\ge2$, $m_{r-1}<m_r$, and an actual fixed-level minimizer,
$f(q_{r-1})\ge0$.
\end{lemma}
\begin{proof}
At $u=q_{r-1}$ the inserted scheme has a zero-variance interval. Raising
its irrelevant mass to the inserted mass and merging the two equal masses
preserves the functional and returns a competitor at the original level.
Fixed-level minimality therefore bounds the actual mass difference quotient
from below by zero. Pass to the checked baseline mass derivative.
The restriction $r\ge2$ preserves the compulsory mass $m_0=0$; this argument
is not claimed for the initial interval.
\end{proof}
The identity for $U''$ is proved below and the full nested second mass bound
is proved above. The following result supplies nonterminal stationarity.

\begin{lemma}[Stationarity at nonterminal levels: Propositions 4.7--4.8]
\lean{SpinGlass.Targets.section4TVarianceQ\_zero\_eq\_overlap\_of\_min}\leanok
\lean{SpinGlass.Targets.hasDerivWithinAt\_section4U\_zero\_of\_min}\leanok
\lean{SpinGlass.Targets.hasDerivWithinAt\_section4FirstVariation\_upper\_zero\_of\_min}\leanok
Suppose the scheme is an actual fixed-level minimizer, $\beta\ne0$,
$m_{r-1}<m_r$, and $r$ is below the final compulsory-mass level
($1\le r\le k$ in the local indexing). Assume
\[
 (q_{r-1}<q_r\ \text{or}\ q_r=0),\qquad
 (q_r<q_{r+1}\ \text{or}\ q_r=1).
\]
Then $Q(m_{r-1},0)=q_r$. Thus the genuine inward derivatives satisfy
$U'(0)=q_r$ and $f'(q_r)=0$. Numerical derivative uniqueness is asserted
only for positive-length intervals.
\end{lemma}
\begin{proof}
The left competitor at upper inserted mass gives $q_r\le Q(0)$ by
Mathlib's endpoint Fermat rule. For the opposite direction, lower the
nonterminal mass $m_r$ to $m_{r-1}$, then insert the original mass at the
next interval. Equal-mass compression gives an actual competitor at the
original level, not merely a formal parameter variation. At its lower
endpoint the actual normalized observable equals the original $Q(0)$,
so Fermat's rule gives $Q(0)\le q_r$. At boundary overlaps use $0\le Q\le1$
for the missing direction. Finally apply the checked inward derivative
identities for $U$ and $f$.
\end{proof}
The next result completes the final compulsory-mass level and handles
coincident interior overlaps. No stationarity assertion is made at $\beta=0$,
where the functional does not constrain overlaps of arbitrary minimizers.

\begin{lemma}[Terminal stationarity and exact reduction]
\lean{SpinGlass.Targets.section4TVarianceQ\_terminal\_zero\_eq\_overlap\_of\_min}\leanok
\lean{SpinGlass.Targets.exists\_strict\_mass\_overlap\_reduction}\leanok
\lean{SpinGlass.Targets.exists\_stationary\_reduction\_of\_min}\leanok
The stationarity conclusion above also holds for $r=k+1$, including $k=0$.
Every original fixed-level minimizer, for $\beta\ne0$, admits an equivalent
scheme of depth $l\le k$ with strict adjacent masses and strict interior
overlaps. It has the same functional, $\psi(t)$, and every $\varphi_N(t)$
on $[0,1]$, remains a minimizer at its reduced depth, and satisfies
$Q(m_{r-1},0)=q_r$ at every physical level. Its first overlap may still be
zero and its last overlap may still be one.
\end{lemma}
\begin{proof}
For terminal stationarity, move the last overlap to one while keeping the
compulsory final mass one. Insert another mass-one level at a variable overlap
and merge the equal masses. This is a genuine fixed-level competitor, whose
endpoint normalized observable is the original $Q(0)$. The same endpoint
Fermat argument supplies the missing inequality.
For reduction, a repeated interior overlap has zero variance, so its mass is
irrelevant to both the scalar and N-site recursions. Raise this mass to the
next one and apply exact equal-mass compression. This preserves strictness
of the remaining mass list. Pad every reduced competitor back to the original
depth to transfer minimality, then iterate. Combine with the earlier mass
reduction. Strict interior overlaps supply both required inward directions;
at zero or one the bound $0\le Q\le1$ supplies the remaining inequality.
\end{proof}
Consequently the conditional quadratic-bound-to-convergence deduction can be
applied only to these reduced schemes and still yield the original unrestricted
Theorem~2.2. The uniform quadratic bound itself is not proved by the reduction.

\begin{lemma}[Actual inserted functional and optimality inputs]
\lean{SpinGlass.Targets.parisiFunctional\_insertLevel}\leanok
\lean{SpinGlass.Targets.section4Phi\_near\_min}\leanok
\lean{SpinGlass.Targets.section4Phi\_upper\_mass\_min}\leanok
The inserted masses and overlaps of (4.27)--(4.29) define an admissible scheme.
Its functional is the actual $\Phi(m,u)$ of (4.37). Near-global minimality and
fixed-level minimality respectively imply
\[
 \begin{aligned}
 \Phi(m,u)&\ge\Phi(m_{r-1},u)-\varepsilon,\\
 \Phi(m_r,u)&\ge\Phi(m_{r-1},u)=\mathcal P_k(m,q).
 \end{aligned}
\]
All mass and overlap endpoints are included.
\end{lemma}
\begin{proof}
Identify every inserted scalar prefix with the actual split cascade. Its
correction follows from the checked paired correction after canceling halved
shared masses. Apply the global infimum to the inserted competitor. At upper
inserted mass, merge the equal adjacent masses to obtain a competitor at the
original level before applying fixed-level minimality.
\end{proof}
These are the inputs (4.30)--(4.31) used in the checked Proposition~4.6 above.

\begin{lemma}[Actual slope velocity and moving squared-slope observable]
\lean{SpinGlass.Targets.hasDerivAt\_stepD1\_parisiF\_variance}\leanok
\lean{SpinGlass.Targets.hasDerivAt\_stepD1\_parisiF\_variance\_curve}\leanok
\lean{SpinGlass.Targets.hasDerivAt\_stepD1\_parisiF\_split\_sq}\leanok
\lean{SpinGlass.Targets.abs\_stepD1Variance\_le\_on\_Icc}\leanok
For an actual Parisi input $A$ and fixed mass $m$, the slope
$B_x(v,x)$ has a genuine positive-variance derivative $J_m(v,x)$.
For $m\ge0$ and $0<v<a$,
\[
 \frac d{dv} B_x(v,x+\sqrt{a-v}\,z)^2
 =2B_x\left(J_m-\frac{z}{2\sqrt{a-v}}B_{xx}\right),
\]
with every term evaluated at the same moving field.
On $0<\ell\le v\le b$ and $0\le m\le1$, the velocity has a common bound
independent of $x$ and the input recursion depth.
\end{lemma}
\begin{proof}
Differentiate the normalized exponential moments defining $B_x$, using
the existing Gaussian-amplitude rule and the checked second-order invariant.
Continuity of both actual partial derivatives gives the joint chain rule.
The tilted Gaussian first-moment estimate bounds $J_m$ uniformly in the field;
the resulting bound is of order $1/\sqrt\ell$ near the variance boundary.
\end{proof}
This supplies the domination prerequisite used in the following cancellation.

\begin{lemma}[Actual negative-square identities (4.16) and (4.45)]
\lean{SpinGlass.Targets.hasDerivAt\_stepD2\_parisiF\_spatial}\leanok
\lean{SpinGlass.Targets.hasDerivAt\_splitBaselineSlopeQ\_parisiF}\leanok
\lean{SpinGlass.Targets.hasDerivAt\_section4TVarianceQ\_baseline}\leanok
\lean{SpinGlass.Targets.hasDerivAt\_deriv\_section4U}\leanok
\lean{SpinGlass.Targets.deriv2\_section4U\_mem\_Icc}\leanok
Let $B(v,x)=T_{m,v}A(x)$ for an actual Parisi input and $0\le m\le1$.
For $0<v<a$, let $\mathbb E_W$ be the actual outer normalized mean with
variance $a-v$ and mass $m$. Then
\[
 \frac d{dv}\mathbb E_W[B_x(v,x+\sqrt{a-v}\,g)^2]
 =-\mathbb E_W[B_{xx}(v,x+\sqrt{a-v}\,g)^2].
\]
After every remaining outer level is retained, at $m=m_{r-1}$ this yields
$Q'(m_{r-1},v)=-R(v)$ with $0\le R(v)\le1$. For $m_{r-1}<1$,
\[
 U''(v)=-R(v)\in[-1,0],\qquad 0<v<\beta^2(q_r-q_{r-1}).
\]
\end{lemma}
\begin{proof}
Differentiate the spatial FTC identity in variance using the common bound
above. The heat equation then gives the genuine third spatial derivative,
without assuming one for the input. The equal-mass semigroup fixes the outer
normalization; dominated differentiation and Gaussian Stein give (4.16).
At baseline all remaining outer potential velocities vanish by (4.11), so
their normalized means propagate the negative-square derivative unchanged.
Finally use the actual $U'=Q$. The $Q'$ identity includes baseline zero and
one; the $U''$ statement inherits the strict upper baseline restriction.
No two-sided variance-endpoint second derivative or strictly negative bound is
claimed. The following result supplies the required inward endpoint calculus.
\end{proof}

\begin{lemma}[Actual Hessian-square continuity and endpoint derivatives]
\lean{SpinGlass.Targets.continuousOn\_section4THessianSquare}\leanok
\lean{SpinGlass.Targets.section4TVarianceQ\_baseline\_sub\_eq\_integral}\leanok
\lean{SpinGlass.Targets.hasDerivWithinAt\_section4TVarianceQ\_baseline}\leanok
\lean{SpinGlass.Targets.derivWithin2\_section4FirstVariation\_eq}\leanok
Let $a=\beta^2(q_r-q_{r-1})$. The actual factor $R$ is continuous on $[0,a]$,
including at baseline mass zero. For $0\le v\le w\le a$,
\[
 Q(m_{r-1},w)-Q(m_{r-1},v)=-\int_v^w R(z)\,dz .
\]
The inward derivative is $Q'=-R$ at both endpoints. For $m_{r-1}<1$ the
actual first variation has inward derivatives
\[
 \begin{aligned}
 f'(u)&=\frac{\beta^2}{2}\bigl(u-Q(m_{r-1},\beta^2(q_r-u))\bigr),\\
 f''(u)&=\frac{\beta^2}{2}\bigl(1-\beta^2R(\beta^2(q_r-u))\bigr).
 \end{aligned}
\]
Numerical derivative uniqueness is asserted only for positive interval length.
\end{lemma}
\begin{proof}
Reuse joint second-observable continuity and the exact baseline semigroup
through every remaining outer mean. Apply the FTC to the existing interior
negative-square derivative, then extend the continuous integrand by interval
projection to obtain inward endpoint derivatives. Apply the within-set chain
rule to the physical variance map for the actual first variation.
\end{proof}

\begin{lemma}[Cubic Taylor and curvature, conditional on actual regularity]
\lean{SpinGlass.Targets.section4FirstVariation\_cubic\_taylor\_of\_hessian\_lipschitz}\leanok
\lean{SpinGlass.Targets.section4FirstVariation\_initial\_short\_curvature\_bound}\leanok
\lean{SpinGlass.Targets.section4FirstVariation\_curvature\_bound\_of\_hessian\_lipschitz\_all\_levels}\leanok
Assume explicitly that the actual $R$ is $L$-Lipschitz on its full physical
variance interval, where $L\ge0$. For $m_{r-1}<1$, put $C=\beta^6L/2$.
For any $u,w\in[q_{r-1},q_r]$, the actual first variation satisfies
\[
 \left|f(u)-f(w)-f'(w)(u-w)-\tfrac12 f''(w)(u-w)^2\right|
 \le C|u-w|^3 .
\]
Suppose in addition that $\beta\ne0$, the scheme is an actual fixed-level
minimizer within $\varepsilon$ of the global infimum, $m_{r-1}<m_r$,
$q_{r-1}<q_r$, and $q_r<q_{r+1}$ or $q_r=1$. At every such physical level,
including the initial and final levels,
\[
 -f''(q_r)\le 2\bigl(C+L_{\mathrm{opt}}(\beta)\bigr)\varepsilon^{1/6},
\]
where $L_{\mathrm{opt}}$ is the checked Proposition~4.6 constant.
\end{lemma}
\begin{proof}
The actual second derivative is $C$-Lipschitz by its displayed formula.
Two applications of Mathlib's endpoint-safe mean-value estimate give the
cubic remainder (with an unoptimized constant). Stationarity makes the
linear term vanish at $q_r$, and $f(q_r)=0$.
For $r\ge2$, test at the lower endpoint when the interval is short, using
$f(q_{r-1})\ge0$, or at distance $\varepsilon^{1/6}$ when it is long,
using Proposition~4.6.
For $r=1$, the compulsory mass $m_0=0$ cannot be varied: no sign for $f(0)$
is assumed. Instead stationarity, $Q\ge0$ and the actual FTC give
\[
 \int_0^{\beta^2q_1}R(v)\,dv\le q_1,
 \qquad -f''(q_1)\le \frac{L\beta^6q_1}{4}.
\]
This handles the short initial interval; the long one uses the same interior
Taylor test. At $\varepsilon=0$ a sufficiently small positive test distance
rules out negative endpoint curvature, without division by the tolerance.
\end{proof}
The regularity hypothesis in this intermediate lemma is discharged by the
following checked results; it is not assumed in the final curvature theorem.

\begin{lemma}[Uniform scalar derivatives and the actual Hessian-square flow]
\lean{SpinGlass.Targets.abs\_parisiFThird\_le\_six}\leanok
\lean{SpinGlass.Targets.abs\_parisiFFourth\_le\_43}\leanok
\lean{SpinGlass.Targets.section4THessianSquare\_lipschitz\_uniform}\leanok
Every actual scalar Parisi input $A$ has genuine spatial derivatives with
$|A'''|\le6$ and $|A''''|\le43$, independently of depth. One further physical
smoothing $B=T_{m,v}A$ has $|B'''|\le14$, $|B''''|\le143$ and
$|\partial_v B''|\le83$ for $v>0$. The actual Section~4 factor satisfies
\[
 |R'(v)|\le535\quad(0<v<a),\qquad
 |R(v)-R(w)|\le535|v-w|\quad(v,w\in[0,a]).
\]
The latter includes both endpoints and degenerate intervals. Zero masses
and variances are retained in the spatial regularity statements.
\end{lemma}
\begin{proof}
The exponential derivative polynomials
\[
 J_3(m)=A'''+3mA'A''+m^2(A')^3,
\]
\[
 J_4(m)=A''''+4mA'A'''+3m(A'')^2+6m^2(A')^2A''+m^3(A')^4
\]
contract under the normalized scalar mean at fixed mass. Their mass-change
bounds are $5\Delta m$ and $42\Delta m$. Starting with $J_3(1)=\tanh$
and $J_4(1)=1$, monotone mass telescoping gives the depth-independent bounds.
The Gaussian heat generator bounds the actual Hessian variance derivative.
For $S=B'$, $T=B''$, $E=B'''$, differentiate the actual initial integral
\[
 R_{\mathrm{initial}}(v)=\mathbb E\bigl[T(x+\sqrt{a-v}\,g)^2W(g)\bigr].
\]
The weight $W$ is the actual equal-mass normalized tilt. Local domination
justifies differentiation. Gaussian Stein applied to
$G=2TE+mT^2S$ removes the inverse outer variance. With bounds $K_2,K_3,K_4$
for $\partial_vT,E,E'$, the result is
$2K_2+K_3^2+K_4+2K_3+2=535$.
All unchanged outer normalized means preserve this bound. Existing endpoint
continuity and the mean-value theorem give the closed-interval estimate.
\end{proof}

\begin{lemma}[Lemma 4.9 and Proposition 4.10 with proved regularity]
\lean{SpinGlass.Targets.abs\_deriv3\_section4FirstVariation\_le\_uniform}\leanok
\lean{SpinGlass.Targets.section4FirstVariation\_cubic\_taylor\_uniform}\leanok
\lean{SpinGlass.Targets.section4FirstVariation\_curvature\_bound\_uniform}\leanok
For $m_{r-1}<1$ and $u\in(q_{r-1},q_r)$ the genuine third derivative satisfies
\[
 f'''(u)=\frac{\beta^6}{2}R'(\beta^2(q_r-u)),\qquad
 |f'''(u)|\le\frac{535\beta^6}{2}.
\]
At $\beta=0$ it is zero. Closed-interval cubic Taylor holds with
$C=535\beta^6/2$. Under the minimality, near-optimality and inward-direction
hypotheses above, Proposition~4.10 holds with that same $C$, without an
assumed derivative or Lipschitz estimate.
\end{lemma}
\begin{proof}
The two signs in the reversed variance coordinate cancel. The genuine
interior second derivative agrees locally with the explicit curvature
coefficient; differentiate it and use the proved bound. Apply the previously
checked endpoint-safe Taylor and short/long-gap curvature argument.
Ordinary derivatives outside the clamped endpoints are not asserted.
\end{proof}

\begin{lemma}[The local-left and initial estimates, Propositions 5.1 and 5.3]
\lean{SpinGlass.Targets.section4THessianSquare\_initial\_le\_zero}\leanok
\lean{SpinGlass.Targets.constrainedPhi\_local\_left\_uniform}\leanok
\lean{SpinGlass.Targets.constrainedPhi\_initial\_left\_uniform}\leanok
Let the actual scheme be a fixed-level minimizer within $\varepsilon$ of
the global infimum, with $\beta\ne0$, $m_{r-1}<m_r$ and the stated right
inward direction. Put
\[
 L_1(\beta)=32+4\cdot535\beta^4+
 \frac{16(535\beta^6/2+L_{\mathrm{opt}}(\beta))}{\beta^2}.
\]
For attainable $u\in[q_{r-1},q_r]$, $0\le t\le t_0<1$, and
\[
 L_1(\beta)\varepsilon^{1/6}\le1-t_0,\qquad
 L_1(\beta)(q_r-u)\le1-t_0,
\]
the actual constrained free energy obeys
\[
 \Phi_N(t,u)\le2\psi(t)-\frac{(1-t_0)^2}{L_1(\beta)}(u-q_r)^2.
\]
For $r=1$, the local-overlap smallness condition can be omitted entirely.
\end{lemma}
\begin{proof}
The actual lambda slope is $Q(t\beta^2(q_r-u))-u$, with inward derivative
$t\beta^2R(t\beta^2(q_r-u))-1$. Stationarity, curvature and the uniform
Lipschitz bound imply a slope at least $(1-t_0)(q_r-u)/2$ locally. Apply
the checked square gain in the actual constrained free energy.
For $r=1$, $m_0=0$: Jensen and the Gaussian semigroup give $R(v)\le R(0)$
throughout $[0,\beta^2q_1]$. The slope bound therefore holds on the entire
initial overlap interval. Coincident endpoints give zero deficit directly.
\end{proof}

\begin{lemma}[The local-right estimate, including the terminal interval]
\lean{SpinGlass.Targets.section5RightV\_baseline\_eq\_reflected\_next}\leanok
\lean{SpinGlass.Targets.section4THessianSquare\_neighbor\_full\_eq\_zero\_all\_levels}\leanok
\lean{SpinGlass.Targets.constrainedPhi\_local\_right\_uniform}\leanok
For $1\le r\le k+1$, put $a=\beta^2(q_{r+1}-q_r)$ and define the actual
reflected factors $Q_+(v)=Q_{r+1}(a-v)$ and $R_+(v)=R_{r+1}(a-v)$,
with baseline mass $m_r$. For $r\le k$, the entire right baseline lambda family
agrees with the next left family at variance $a-v$; at the terminal level the
same identification is obtained on the exactly padded scheme. Moreover,
\[
 Q_+'(v)=R_+(v),\qquad
 |R_+(v)-R_+(w)|\le535|v-w|,
\]
with inward endpoint derivatives, and
\[
 Q_+(0)=Q_r(0),\qquad R_+(0)=R_r(0).
\]
Under the same fixed-level minimality, near-optimality, strict mass and
inward-direction hypotheses as the left estimate, with $q_{r-1}<q_r$,
the local-right estimate holds:
\[
 \begin{gathered}
 q_r\le u\le q_{r+1},\quad 0\le t\le t_0<1,\\
 L_1(\beta)\varepsilon^{1/6}\le1-t_0,\quad
 L_1(\beta)(u-q_r)\le1-t_0\\
 \Longrightarrow\quad
 \Phi_N(t,u)\le2\psi(t)-\frac{(1-t_0)^2}{L_1(\beta)}(u-q_r)^2.
 \end{gathered}
\]
Here $\beta\ne0$, $u$ is attainable, and $L_1$ is the same constant as above.
\end{lemma}
\begin{proof}
For $r\le k$, the paired mass arrays and sharing cutoffs agree exactly; only the split
variance is reversed. Reuse the checked left lambda derivative and factor
calculus. At full inner variance the scalar derivative recursion is the
original next scalar derivative; the zero outer variance evaluates directly.
This proves the neighboring endpoint identities through all remaining means.
For $r=k+1$, append a redundant mass-one level at overlap one. The added
innermost step has zero variance. Exact identities shift the scalar recursion
and its first two derivatives, the full paired lambda family, and the normalized
$Q/R$ recursions by one level. Thus the same factor calculus and endpoint
identities apply to the original terminal interval. Only the original scheme's
minimality is used; minimality of the padded scheme is not asserted.
Stationarity and curvature therefore transfer from the original level.
The right slope $Q_+(t\beta^2(u-q_r))-u$ has derivative
$t\beta^2R_+(t\beta^2(u-q_r))-1$. Locally it is at most
$-(1-t_0)/2$, so the slope is at most $-(1-t_0)(u-q_r)/2$.
Apply the checked right lambda square gain for the actual constrained free energy.
\end{proof}

\begin{lemma}[Zero-first-overlap curvature and the completed local-right case]
\lean{SpinGlass.Targets.hasDerivWithinAt\_zeroOuterSquaredSlope\_zero}\leanok
\lean{SpinGlass.Targets.section4THessianSquare\_initial\_zero\_le\_of\_min}\leanok
\lean{SpinGlass.Targets.constrainedPhi\_local\_right\_of\_reduced\_min}\leanok
Suppose $q_1=0$, $q_2>0$, $m_1>0$, $\beta\ne0$, and the original scheme
minimizes the functional at its fixed level. Then
\[
 \beta^2R_1(0)\le1.
\]
The local-right quadratic estimate therefore holds at $r=1$ with zero first
overlap, with the same constant $L_1(\beta)$ and without a near-optimality
assumption in this boundary case. Combining with the positive-left-gap
estimate proves Proposition~5.2 for every physical level of a reduced scheme
in the exact-covariance SK setting.
\end{lemma}
\begin{proof}
Let $A$ be the original scalar input immediately above the first mass,
$m=m_1$, and $a=\beta^2q_2$. Write $B_v$ for the scalar Parisi step of
mass $m$ and variance $v$ applied to $A$, and set
\[
 G(u)=\mathbb E\bigl[B_{a-u}'(h+\sqrt{u}Z)^2\bigr],\qquad Z\sim N(0,1).
\]
The checked scalar-square identities and stationarity give $B_a'(h)=0$.
For each fixed $z$, positive-variance calculus applied to
$B_{a-\theta^2}'(h+\theta z)$ at $\theta=0$ gives derivative $B_a''(h)z$.
After division by $\theta$ and squaring, a quadratic Gaussian envelope
permits dominated convergence. Thus the genuine inward derivative satisfies
\[
 G(0)=0,\qquad G_+'(0)=B_a''(h)^2=R_1(0).
\]
The existing auxiliary-base constructions give an actual same-level
competitor $f$ minimized at zero, with closed-interval derivative
\[
 f'(u)=\frac{\beta^2m_1}{2}\bigl(u-G(\beta^2u)\bigr).
\]
Its outer mass is zero and its inner mass is $m_1$; no equal-mass derivative
formula is substituted. The first derivative vanishes at zero, and
\[
 f_+''(0)=\frac{\beta^2m_1}{2}\bigl(1-\beta^2R_1(0)\bigr)\ge0.
\]
The last inequality is the one-sided second-order necessary condition,
proved using the slope limit and mean-value theorem. Dividing by the positive
coefficient gives the curvature bound. Apply the already checked right-factor
calculus and lambda gain with zero curvature error. When $q_2=0$, the right
interval is a singleton and the non-strict pressure bound suffices.
The terminal replica-symmetric case uses the existing terminal auxiliary base.
\end{proof}
The outside-neighbor cases, compactness, and the uniform assembly into
Theorem~2.4 are still open; the far-right and signed initial cases are proved below.

\begin{lemma}[The actual right mass variation and its optimality bound]
\lean{SpinGlass.Targets.parisiFunctional\_insertRightLevel}\leanok
\lean{SpinGlass.Targets.hasDerivAt\_section4RightT\_mass\_baseline}\leanok
\lean{SpinGlass.Targets.section4RightT\_mass\_taylor\_bound}\leanok
\lean{SpinGlass.Targets.section4RightFirstVariation\_upper\_bound}\leanok
Let $1\le r\le k+1$, $a=\beta^2(q_{r+1}-q_r)$, and let $T_+(m,v)$
be the actual right scalar cascade. Its split is
\[
 P_{m,v}\bigl(P_{m_r,a-v}A\bigr),
\]
followed by the unchanged outer levels, where $P$ is the scalar Parisi step.
The inserted outer mass varies; the original inner mass remains fixed.
The genuine derivative exists on a two-sided neighborhood of every mass in
$[0,1]$, including both physical variance endpoints. Define
\[
 U_+(v)=2\partial_mT_+(m_r,v),\qquad
 f_+(u)=\frac12U_+\bigl(\beta^2(u-q_r)\bigr)
       -\frac{\beta^2}{4}(u^2-q_r^2).
\]
For an original fixed-level minimizer within $\varepsilon$ of the infimum,
with $m_{r-1}<m_r$ and $q_r\le u\le q_{r+1}$,
\[
 f_+(u)\le O_\beta\sqrt\varepsilon.
\]
Here $O_\beta$ is the same depth-independent constant used in Proposition~4.6.
Zero lower mass and the terminal mass $m_r=1$ are included.
\end{lemma}
\begin{proof}
Insert the overlap $u$ and a mass $m\in[m_{r-1},m_r]$ into the original
scheme. Its actual functional is
\[
 \log2+T_+\bigl(m,\beta^2(u-q_r)\bigr)
 -\operatorname{corr}(s)
 +(m_r-m)\frac{\beta^2}{4}(u^2-q_r^2).
\]
At $m=m_r$ this is the original functional. At $m=m_{r-1}$, merge the
adjacent equal masses to obtain a same-level competitor; original minimality
applies without assuming minimality of a padded scheme. Near-global
minimality applies to every intermediate insertion. Reuse the scalar mass
derivatives and the normalized second-derivative covariance invariant through
the fixed outer cascade. This yields the same uniform quadratic Taylor
remainder as on the left. Apply the existing quadratic comparison lemma to
$-f_+$, since the admissible scalar variation decreases the mass.
\end{proof}

\begin{lemma}[Right pressure improvement from a negative mass derivative]
\lean{SpinGlass.Targets.constrainedPhi\_le\_two\_guerraPsi\_right\_mass\_variation}\leanok
\lean{SpinGlass.Targets.exists\_constrainedPhi\_right\_mass\_gap\_uniform\_in\_size}\leanok
For $m_r>0$, $0\le t\le1$, and $q_r\le u\le q_{r+1}$, set
\[
 v=t\beta^2(u-q_r),\qquad
 D_+=U_+(v)-t\frac{\beta^2}{2}(u^2-q_r^2).
\]
If $D_+<0$, there is $c>0$, chosen before $N$ and the disorder, such that
\[
 \Phi_N(t,u)\le2\psi(t)-c
\]
for every positive system size at which $u$ is attainable.
\end{lemma}
\begin{proof}
The actual paired interpolation admits scalar masses up to $2m_r$;
its inserted paired mass is $m/2$. The corrected scalar trial function has
derivative $D_+$ at $m_r$, so an admissible \emph{increase} in mass gives a
strict improvement, even at $m_r=1$. Transport that scalar deficit using the
proved right endpoint bound and zero-lambda identity. If the lambda slope is
nonzero, the existing lambda square gain already gives a positive deficit.
The mass-or-lambda dichotomy reduces the remaining case to the actual scalar
derivative sign proved below.
\end{proof}

\begin{lemma}[Right supporting line and the far-right case, Proposition 5.6]
\lean{SpinGlass.Targets.section4RightU\_sub\_eq\_integral}\leanok
\lean{SpinGlass.Targets.hasDerivWithinAt\_section4RightU}\leanok
\lean{SpinGlass.Targets.convexOn\_section4RightU}\leanok
\lean{SpinGlass.Targets.exists\_constrainedPhi\_far\_right\_uniform\_in\_size}\leanok
For $m_r>0$, the actual right mass derivative satisfies
\[
 U_+(w)-U_+(v)=\int_v^w Q_+(z)\,dz
 \quad(0\le v\le w\le a),
\]
so $U_+'=Q_+$, including inward derivatives at both endpoints.
Since $Q_+'=R_+\ge0$, $U_+$ is convex. For
$v=\beta^2(u-q_r)$ this yields
\[
 \begin{split}
 D_+\le{}&2f_+(u)-(1-t)\frac{\beta^2}{2}(u-q_r)^2\\
 &-(1-t)v\bigl(Q_+(tv)-u\bigr).
 \end{split}
\]
Consequently Proposition~5.6 holds for every physical interval of a reduced
scheme, with $L_3(\beta)=4O_\beta$. Specifically, for $\beta\ne0$, $L_1>0$,
$0\le t\le t_0<1$, original fixed-level minimality and $\varepsilon$-near
global minimality, the conditions
\[
 \begin{gathered}
 q_r\le u\le q_{r+1},\quad L_1(u-q_r)\ge1-t_0,\\
 L_3(\beta)\sqrt\varepsilon\le
 (1-t_0)\frac{\beta^2}{2}\left(\frac{1-t_0}{L_1}\right)^2
 \end{gathered}
\]
give a strictly positive constrained-free-energy deficit chosen before $N$
and the disorder. This includes $r=k+1$ and $m_r=1$.
\end{lemma}
\begin{proof}
Remove the mass-gap prefactor from the actual right variance derivative to
obtain a bounded normalized observable. At baseline, identify its full
recursion with the existing reflected $Q_+$, not just their products with a
zero mass gap. Mass continuity and bounded dominated convergence justify a
left mass difference quotient of the variance integral identity. The positive
baseline mass permits this limit even at mass one. A continuous extension of
$Q_+$ proves the inward endpoint derivative by the fundamental theorem of
calculus. Apply the convex supporting line and the proved upper bound on
$f_+$. If $Q_+(tv)=u$, the smallness condition makes $D_+<0$; otherwise the
lambda square gain is strict. Apply the preceding mass-or-lambda result.
No compactness over time and overlap is claimed in this proposition.
\end{proof}

\begin{lemma}[Signed initial Gaussian calculus]
\lean{SpinGlass.Targets.hasDerivAt\_signedSplitSlope}\leanok
\lean{SpinGlass.Targets.signedSplitHessian\_le\_unsplit}\leanok
\lean{SpinGlass.Targets.signedSplitSlope\_parisiF\_lower\_bound}\leanok
Let $A$ be the actual scalar cascade immediately above the initial zero mass,
$a=\beta^2q_1$, and $B_v=P_{0,v}A$. Define
\[
 W(v)=\mathbb E\bigl[B_v'(h+\sqrt{a-v}Z)B_v'(h-\sqrt{a-v}Z)\bigr].
\]
The function $W$ is continuous on $[0,a]$, $W(a)=B_a'(h)^2\ge0$, and
for $0<v<a$ its genuine derivative satisfies
\[
 W'(v)=\mathbb E\bigl[B_v''(h+\sqrt{a-v}Z)B_v''(h-\sqrt{a-v}Z)\bigr]
 \le P_{0,a}\bigl((A'')^2\bigr)(h)=R_1(0).
\]
Consequently $W(v)\ge-R_1(0)(a-v)$, including both endpoints.
\end{lemma}
\begin{proof}
Differentiate the two actual smoothed slopes using the existing scalar
variance and spatial calculus. Gaussian integration by parts cancels the
third-derivative terms. The pointwise inequality $2xy\le x^2+y^2$,
Gaussian reflection, and the checked zero-mass Jensen/semigroup bound give
the derivative estimate. Apply the closed-interval mean-value inequality
and the nonnegative independent endpoint. No new regularity assumption or
stationarity identity is introduced.
\end{proof}

\begin{lemma}[Signed scalar lambda families and the retained-field endpoints]
\lean{SpinGlass.Targets.section5SignedInitialV\_zero}\leanok
\lean{SpinGlass.Targets.hasDerivAt\_section5SignedInitialV\_zero}\leanok
\lean{SpinGlass.Targets.section5RetainedSignedInitialV\_zero}\leanok
\lean{SpinGlass.Targets.section5RetainedSignedInitialV\_lambda\_gain}\leanok
\lean{SpinGlass.Targets.section5SignedInitialInterpolation\_one}\leanok
\lean{SpinGlass.Targets.section5SignedInitialInterpolation\_zero}\leanok
The purely signed scalar family has zero-lambda value $2A_0(h)$ and
lambda derivative $W(v)$ at zero. Its second lambda derivative lies in
$[0,1]$, giving the usual lambda square gain.

The interpolation to the original constrained free energy must also retain
the frozen \emph{positive} shared field. For $-q_1\le u\le0$ set
\[
 v=t\beta^2(q_1+u),\qquad b=(1-t)\beta^2q_1,\qquad
 c=-t\beta^2u.
\]
These variances are nonnegative and $v+b+c=a$. The retained scalar family
uses independent variance $v$, positive shared variance $b$, and negative
shared variance $c$. Its zero-lambda value is still $2A_0(h)$, its actual
lambda derivative is the corresponding two-Gaussian slope product, and
its second lambda derivative again lies in $[0,1]$.

The checked signed interpolation has exactly the original constrained
free energy at time one and the constrained retained-field endpoint at time
zero. Finite-site tensorization bounds the latter by
$2\log2+V_{\mathrm{ret}}(\lambda)-\lambda u$; it is not identified with
the unconstrained scalar family itself. The original external field is unchanged.
\end{lemma}
\begin{proof}
Apply the existing finite-step lambda curvature invariant to the independent
prefix and the two mass-zero signed Gaussian steps. Each marginal law is
unchanged, so the scalar semigroup proves the zero-lambda identity.
At interpolation time one the added negative variance is zero; the existing
exact cascade reindexing recovers the original constrained free energy.
At time zero the disorder vanishes and the variance arrays give the stated
field endpoint.
\end{proof}
\begin{lemma}[Quantitative retained-field zero-time gain]
\lean{SpinGlass.Targets.section5RetainedSignedInitialV\_deriv\_lower\_bound}\leanok
\lean{SpinGlass.Targets.section5SignedInitialFieldEndpoint\_quantitative\_gain}\leanok
For the variances $v,b,c$ above, the actual retained-family derivative obeys
\[
 \partial_\lambda V_{\mathrm{ret}}(0)\ge-cR_1(0).
\]
If $0\le t\le t_0<1$, $\beta^2R_1(0)\le1+e$, and
$0\le e\le(1-t_0)/2$, then its lambda slope is at least
$(1-t_0)(-u)/2$. Hence, for positive $N$ and attainable overlap, the genuine
zero-time constrained field endpoint is at most
\[
 2\log2+2A_0(h)-\frac{(1-t_0)^2}{8}u^2.
\]
\end{lemma}
\begin{proof}
Condition on the positive frozen Gaussian. Apply the generic signed estimate
with total variance $v+c$ and shifted external field. Bounded Fubini and the
scalar semigroup combine the remaining Hessian-square mean with variance $b$
to give $R_1(0)$. The curvature assumption bounds the slope; apply the
unit-curvature lambda gain and the checked zero-time endpoint comparison.
\end{proof}
The endpoint gain is transported to the original pressure by the following
exact reflection and conditioning argument. The retained positive shared
field is never replaced by a negative one.

\begin{lemma}[Exact signed initial pressure transport]
\lean{SpinGlass.Targets.skDisorder\_even\_ae}\leanok
\lean{SpinGlass.Targets.section5SignedInitialInterpolation\_eq\_conditioned\_average}\leanok
\lean{SpinGlass.Targets.section5SignedInitialInterpolation\_endpoint\_bound}\leanok
For $N>0$, $0\le t\le1$, and attainable $-q_1\le u\le0$, the actual
signed interpolation satisfies
\[
 \Psi(t,u)\le \mathcal E_0(t,u)-2tC,
\]
where $\mathcal E_0$ is the retained-field zero-time endpoint above and $C$
is the Parisi correction in $\psi(t)=\log2+A_0(h)-tC$.
\end{lemma}
\begin{proof}
Exact SK covariance gives
$\E[(U(-\sigma)-U(\sigma))^2]=0$, hence Hamiltonian evenness on a common
full-measure set. The formal proof uses the existing spectral representation
and includes zero spectral variances. Reflecting the second replica reverses
its overlap and physical field. Reflection of its independent Gaussian fields
also preserves their product law.

Condition on the frozen positive field $\sqrt b\,z$, with
$b=(1-t)\beta^2q_1$. The two remaining external fields are explicitly
$x=h+\sqrt b\,z$ and $-x$. The negative shared interpolating variance becomes
the ordinary positive shared variance $(1-w)t\beta^2(-u)$.
The existing arbitrary-field trace and heat calculus, with the frozen variance
removed from this conditional path, gives its genuine derivative bound
$\partial_w\widetilde\Psi(w;x,-x)\le-2tC$ and closed-interval endpoint
inequality. Joint continuity and affine Gaussian bounds justify Fubini and
integration over $z$. Undoing reflection recovers exactly the original
constrained free energy, including $t=0$ and all zero-variance endpoints.
This proves the required signed transport without asserting a general signed
version of Theorem~3.1.
\end{proof}

\begin{proposition}[The negative initial interval: Proposition 5.4]
\lean{SpinGlass.Targets.section5Initial\_curvature\_data}\leanok
\lean{SpinGlass.Targets.constrainedPhi\_initial\_signed\_uniform}\leanok
\lean{SpinGlass.Targets.constrainedPhi\_initial\_signed\_lt}\leanok
Let $\beta\ne0$ and let the scheme minimize at its own fixed level and be
within $\varepsilon$ of the infimum of the Parisi functional. Assume the
reduced-scheme conditions $m_0<m_1$ and $q_1<q_2$ or $q_1=1$.
With the same explicit beta-only constant $L_2$ as in Proposition~5.3,
assume $L_2\varepsilon^{1/6}\le1-t_0$ and $0\le t\le t_0<1$.
For $N>0$ and every attainable $-q_1\le u<0$,
\[
 \Psi(t,u)\le2\psi(t)-\frac{(1-t_0)^2}{8}u^2<2\psi(t).
\]
\end{proposition}
\begin{proof}
The initial curvature argument already used in Proposition~5.3 supplies
$0\le e\le(1-t_0)/2$ with $\beta^2R_1(0)\le1+e$; its regularity input is
the checked universal Hessian-square Lipschitz constant $535$.
Apply the retained-field zero-time gain and the exact pressure transport.
The deficit is strictly positive because $u<0$ and $t_0<1$.
If $q_1=0$ the negative interval is empty. The non-strict curvature-based
estimate also includes $u=0$, but the displayed deficit vanishes there;
the uniform assembly of Theorem~2.4 is a separate remaining obligation.
\end{proof}

\begin{lemma}[Tagged mass interleaving for Proposition 5.7]
\lean{SpinGlass.Targets.monotone\_section5InterleavedMass}\leanok
\lean{SpinGlass.Targets.section5Interleaving\_correction}\leanok
\lean{SpinGlass.Targets.section5Interleaving\_variance\_grouping}\leanok
Use the paper's indices. The physical list contains $k+1$ masses and the
inserted interpolation list contains $k+2$. Their tagged union admits a
stable nondecreasing ordering with endpoints zero and one. The source
embeddings preserve order, including all equal-mass ties. Attaching the
original frozen variances to physical tags and the inserted variances to
interpolation tags gives nonnegative variances for $0\le t\le1$ and
$q_{s-1}\le |u|\le q_s$.

The correction sum is exactly (5.36). After doubling the scalar mass of a
shared tag, every resulting mass is an original $m_p$. For every weight
function $W$ on original levels, the assigned variances satisfy
\[
 \sum_\ell W(p_\ell)a_\ell
   =\sum_{p=0}^{k}W(p)\beta^2(q_{p+1}-q_p).
\]
Thus (5.49) also holds with repeated numerical masses, by grouping all their
original levels together. The local Lean scheme has paper depth $k+1$, so
its tagged union has $2k+5$ entries.
\end{lemma}
\begin{proof}
Use Mathlib's stable tuple sort on the concatenated lists, retaining source
and index tags. Exact finite-sum transport through its permutation reduces
the identities to the already checked insertion correction and telescoping
variance sums. No tag is identified with another merely because their
numerical masses agree.
\end{proof}

\begin{lemma}[Scalar operator interchange and sorting: (5.46)]
\lean{SpinGlass.Targets.parisiStep\_interchange}\leanok
\lean{SpinGlass.Targets.parisiListCascade\_insertionSort\_le}\leanok
\lean{SpinGlass.Targets.parisiListCascade\_merge}\leanok
Let $Q$ be measurable with linear growth, $a,a'\ge0$, and $0\le m'\le m$.
Then the actual Gaussian Parisi operators satisfy
\[
 T_{m,a}T_{m',a'}Q\le T_{m',a'}T_{m,a}Q.
\]
Consequently, sorting a finite outermost-first list of nonnegative-mass,
nonnegative-variance operators by increasing mass raises its composition.
Adjacent equal-mass operators merge their variances exactly, even at mass zero.
\end{lemma}
\begin{proof}
For $0<m'<m$, integral Minkowski follows from Mathlib's H\"older inequality
and Tonelli. Linear growth ensures the required exponential integrals are
finite and identifies them with the nested log means. Analytic continuation
and monotonicity in the lower mass give the case $m'=0$; the Gaussian
semigroup handles equal masses. Induction over insertion sort repeatedly
applies this actual interchange and scalar order preservation.
\end{proof}
Nonconstancy alone is insufficient for a strict comparison: affine inputs
also give equality. The checked strict result below uses the actual positive
scalar Hessian and includes the zero lower-mass face.

\begin{lemma}[Strict curvature and Gaussian equality rigidity]
\lean{SpinGlass.Targets.scalarFieldCascade\_C2\_pos}\leanok
\lean{SpinGlass.Targets.strictConvexOn\_scalarFieldCascade}\leanok
\lean{SpinGlass.Targets.scalarFieldCascade\_even}\leanok
\lean{SpinGlass.Targets.eq\_mul\_of\_gaussian\_cauchySchwarz\_eq}\leanok
Every finite scalar recursion from $\log\cosh$ with masses in $[0,1]$ has a
strictly positive actual Hessian, a strictly increasing slope, and is even
and strictly convex. The mass order need not be sorted, and zero variances
are included. Moreover, equality in Gaussian Cauchy--Schwarz for continuous
profiles $f,g$, with $g>0$ and the three required moments integrable, implies
$f=cg$ everywhere for some constant $c$.
\end{lemma}
\begin{proof}
Positive input curvature gives a positive tilted Hessian mean; the additional
mass-weighted slope variance is nonnegative. Induct over the actual scalar
recursion and apply the second-derivative criterion for strict convexity.
Gaussian reflection preserves evenness. For equality rigidity, integrate
$(f-cg)^2$ with $c=\E[fg]/\E[g^2]$. Equality makes this integral zero.
Gaussian absolute continuity transfers the resulting almost-everywhere
identity to Lebesgue measure; continuity makes it pointwise.
\end{proof}
\begin{lemma}[One-step paired comparison: Lemma 5.10 on actual scalar inputs]
\lean{SpinGlass.Targets.independentStepPi\_one\_le\_scalar}\leanok
\lean{SpinGlass.Targets.gtScalarStep\_shared\_le}\leanok
\lean{SpinGlass.Targets.gtScalarStep\_opposite\_le}\leanok
\lean{SpinGlass.Targets.independentStepPi\_one\_lt\_scalar\_of\_off\_diagonals}\leanok
\lean{SpinGlass.Targets.gtScalarStep\_shared\_lt\_of\_strictMono\_deriv}\leanok
\lean{SpinGlass.Targets.gtScalarStep\_opposite\_lt\_of\_strictMono\_deriv}\leanok
Suppose $D(x,y)\le G(x)+G(y)$, with measurable linear-growth inputs and
the corresponding Gaussian integrability. At raw paired mass $m\ge0$,
the independent step is bounded by $T_{m,a}G(x)+T_{m,a}G(y)$, and the shared
and opposite-shared steps by $T_{2m,a}G(x)+T_{2m,a}G(y)$.

For the genuine differentiable scalar inputs with strictly increasing slope,
if $m,a>0$, shared comparison is strict when $x\ne y$. With even $G$,
opposite comparison is strict when $x\ne-y$. If the input comparison is
strict off $x=\pm y$ and $a>0$, the independent step makes it strict
everywhere, including at mass zero. An almost-everywhere strict comparison
also propagates through any signed or independent Gaussian step.
\end{lemma}
\begin{proof}
Use independence for product factorization and Gaussian Cauchy--Schwarz for
the shared cases. Equality forces proportional translated exponential
profiles. Their logarithmic difference is constant, so the actual strictly
increasing derivative forces equal shifts. Reflection gives the opposite
case. For independent propagation, non-atomicity makes both exceptional
diagonals null; strict integral monotonicity then applies. All exponential
integrability follows from the stated growth bounds.
\end{proof}
\begin{lemma}[Strict scalar interchange: (5.47) on the actual inputs]
\lean{SpinGlass.Targets.tilt\_variance\_pos\_of\_strictMono}\leanok
\lean{SpinGlass.Targets.parisiStep\_interchange\_strict}\leanok
\lean{SpinGlass.Targets.parisiStep\_interchange\_parisiF\_strict}\leanok
For $0\le m'<m\le1$, $a,b>0$, and the actual scalar inputs $A$,
\[
 T_{m,a}T_{m',b}A(x)<T_{m',b}T_{m,a}A(x).
\]
More generally, the formal statement takes linear growth, the established
bounded C2 data, and a continuous strictly positive Hessian.
\end{lemma}
\begin{proof}
Differentiate the genuine commutator
$T_{m,v}T_{m',b}A-T_{m',b}T_{m,v}A$ inward at $v=0$.
The derivative is $-(m-m')/2$ times the tilted variance of $A'$.
That variance is strictly positive: equality in Cauchy--Schwarz would make
the continuous slope constant under a full-support Gaussian, contrary to
its strict increase. A small strict swap therefore exists. Split any
positive upper variance by the same-mass semigroup; strict order through
the remaining positive-variance outer step and non-strict interchange
give the displayed inequality. This argument applies also when $m'=0$.
\end{proof}

\begin{lemma}[Actual sorted scalar endpoint and its equality order]
\lean{SpinGlass.Targets.section5SortedScalarSteps\_cascade\_eq\_parisiF}\leanok
\lean{SpinGlass.Targets.section5InterleavedScalarSteps\_cascade\_le\_parisiF}\leanok
\lean{SpinGlass.Targets.section5InterleavedScalarSteps\_eq\_implies\_order}\leanok
Sorting the actual interleaved scalar operators yields exactly the original
Parisi recursion $A_0$. Hence $D_0^*(h)\le A_0(h)$, as in (5.50).
Equality at one field forces the scalar masses to be nondecreasing on
positive-variance tags, as in (5.51). Zero-variance tags impose no order.
\end{lemma}
\begin{proof}
Equal-mass operators commute and merge by the Gaussian semigroup. Group
the sorted tags by original level and use (5.49). This includes accidental
mass ties. For strictness, insertion-sort induction locates an inverted
pair of positive-variance operators and applies strict interchange to its
actual log-cosh suffix. Strict order propagates through all outer operators;
a zero-variance outer operator is exactly the identity.
\end{proof}

\begin{proposition}[Mixed paired comparison and the order in Proposition 5.11]
\lean{SpinGlass.Targets.mixedScalarCascade\_zero\_le}\leanok
\lean{SpinGlass.Targets.mixedScalarCascade\_eq\_implies\_order}\leanok
\lean{SpinGlass.Targets.mixedScalarCascade\_lt\_of\_opposite\_before\_shared}\leanok
The genuine mixed paired recursion satisfies
$D_0(x,y)\le D_0^*(x)+D_0^*(y)$ for arbitrary order of independent, shared
and opposite-shared steps. Equality at $(h,h)$ forces each positive-variance
shared tag with positive mass to precede every positive-variance independent
tag in the paper's forward indexing, as in (5.44).
An active opposite-shared step followed outward by an active ordinary-shared
step also forces strict comparison everywhere.
\end{proposition}
\begin{proof}
Iterate the checked one-step comparisons, retaining the actual Gaussian
integrals and RSAT's finite-step regularity. A positive-mass shared step
creates strictness off its diagonal. It persists off both diagonals through
intervening signed steps; a positive-variance independent step spreads it
everywhere by Gaussian non-atomicity. For the opposite/shared obstruction,
propagate strictness off the anti-diagonal until a shared step spreads it.
All zero-mass and zero-variance branches are treated explicitly.
\end{proof}

\begin{lemma}[Cumulative overlaps and the original interleaved endpoint]
\lean{SpinGlass.Targets.section5Interleaved\_pairCorrection}\leanok
\lean{SpinGlass.Targets.section5InterleavedInterpolation\_one}\leanok
\lean{SpinGlass.Targets.section5InterleavedInterpolation\_zero\_le\_lambda}\leanok
The cumulative overlap at merged position $\ell$ is the inserted overlap
indexed by the number of interpolation tags before $\ell$. It advances
only at interpolation tags, has endpoints zero and one, and equals $|u|$
at the distinguished cutoff. Its actual signed correction is (5.36).
The full interleaved interpolation satisfies $\eta(1)=\Psi(t,u)$, for
either sign of $u$, including zero variances and physical-time endpoints.
For every attainable $u$ and every $\lambda\in\R$, its other endpoint obeys
$\eta(0)\le2\log2+D_0(\lambda;h,h)-\lambda u$, with the actual mixed
paired scalar recursion and the selected interval containing $|u|$.
\end{lemma}
\begin{proof}
Strict source-order embeddings identify the cumulative increments and
recover the physical list after filtering interpolation tags. At second
time one, every interpolation variance is zero, so its actual operator is
the identity, including at zero mass. Gaussian Fubini identifies the
remaining independent operators, and the checked physical-field rescaling
recovers exactly the original constrained free energy. Physical sharing
remains positive even when the interpolating sharing is opposite.
At second time zero the disorder vanishes exactly. The genuine constrained
terminal relaxation passes through the mixed operators by Gaussian order;
finite-site tensorization supplies the scalar expression. Exact index
reversal identifies the forward tagged list with the bottom-up recursion.
\end{proof}

\begin{lemma}[The separate time-zero deficit]
\lean{SpinGlass.Targets.constrainedPhi\_time\_zero\_le\_guerraPsi\_of\_min}\leanok
For every attainable $u$, at each reduced minimizing level,
\[
 \Psi(0,u)\le2\psi(0)-\frac{(u-q_r)^2}{2}.
\]
The strict conclusion holds whenever $u\ne q_r$.
\end{lemma}
\begin{proof}
At physical time zero the second interpolation is constant. Apply the
existing arbitrary-overlap constrained endpoint comparison and quadratic
lambda gain. Its derivative is the actual $Q_r(0)$; all-level stationarity
identifies this with $q_r$. No neighboring-overlap restriction is needed.
\end{proof}

\begin{lemma}[Strict actual field endpoint in the outside-index regimes]
\lean{SpinGlass.Targets.section5InterleavedScalarV\_zero\_lt\_left\_outside}\leanok
\lean{SpinGlass.Targets.section5InterleavedScalarV\_zero\_lt\_right\_outside}\leanok
\lean{SpinGlass.Targets.exists\_section5Interleaved\_zero\_gap\_left\_outside}\leanok
\lean{SpinGlass.Targets.exists\_section5Interleaved\_zero\_gap\_right\_outside}\leanok
Assume strict original masses and choose the interval containing $|u|$.
When its index is outside the two neighboring indices and the two tagged
witness variances are positive, the actual recursions satisfy
$D_0(h,h)<2A_0(h)$. Consequently, $\delta=2A_0(h)-D_0(h,h)>0$ is independent
of system size and disorder, and
\[
 \eta(0)\le2\log2+2A_0(h)-\delta.
\]
Both signs of $u$ are allowed. The formal statements retain the selected
interval range and positive witness variances explicitly.
\end{lemma}
\begin{proof}
Equality in the joint comparison would give equality in both the paired
and scalar comparisons. The proved implications yield (5.44) and (5.51),
contradicting the checked outside-neighbor tagged order obstruction.
Transport the resulting scalar deficit through the genuine zero endpoint
bound. The deficit is fixed before quantifying over $N$ and the disorder.
\end{proof}

\begin{lemma}[Actual mixed calculus and endpoint continuity]
\lean{SpinGlass.Targets.hasDerivAt\_mixedConstrainedCascade}\leanok
\lean{SpinGlass.Targets.hasDerivAt\_mixedConstrainedCascade\_second}\leanok
\lean{SpinGlass.Targets.hasDerivAt\_mixedConstrainedCascade\_variance}\leanok
\lean{SpinGlass.Targets.continuousOn\_section5InterleavedInterpolation}\leanok
\lean{SpinGlass.Targets.differentiableAt\_section5InterleavedIntegrand}\leanok
For arbitrary ordering of independent, shared and opposite-shared Gaussian
levels, the genuine constrained recursion has the normalized-mean first
derivatives and covariance second derivatives. The disorder Hessian is bounded
by $(8+16\sum_\ell |m_\ell|)\|V\|_1\|W\|_1$; the spatial analogue is bounded
by $(2+4\sum_\ell |m_\ell|)(\|A\|_1+\|B\|_1)(\|C\|_1+\|D\|_1)$.
The derivative in one positive variance is its Gaussian heat generator
transported through all remaining mixed levels. Every unchanged variance
and every mass may be zero. The actual interleaved free energy $\eta(w)$ is
continuous for $0\le w\le1$, for either overlap sign and including zero faces.
Its actual tagged integrand is differentiable for $0<w<1$, also at physical
time zero. This follows from genuine joint Fr\'echet differentiability in
time and both fields; the explicit disorder-plus-heat decomposition of this
derivative remains a separate obligation.
\end{lemma}
\begin{proof}
Reuse the established one-step parameter and covariance rules. Opposite
sharing is a fixed second-field reflection. The independent two-step
covariance has the common form by cancellation of the intermediate products.
The existing finite-dimensional Gaussian heat theorem applies to these actual
derivatives. Compact-parameter continuity passes through each mixed Gaussian
operator; contraction in the terminal disorder supplies domination of the
outer Gaussian expectation.
\end{proof}

\begin{lemma}[Actual mixed replica and disorder identities]
\lean{SpinGlass.Targets.mixedConstrainedCascadeDD\_disorder\_eq\_replica}\leanok
\lean{SpinGlass.Targets.mixedOuterMean\_spatialHeat\_eq\_replica}\leanok
\lean{SpinGlass.Targets.hasDerivAt\_mixedConstrainedGaussian\_amplitude\_trace}\leanok
\lean{SpinGlass.Targets.mixedConstrainedSecond\_SK\_trace}\leanok
The genuine mixed Gibbs weights and product-before-transport split weights
are normalized and nonnegative. The actual Hessian has the finite covariance
telescope under those weights, and the spatial Hessian-plus-square term
retains its original split after any number of outer transports. With field
variances fixed, differentiation of the actual outer Gaussian average in the
disorder amplitude equals that amplitude times the expected spectral Hessian
trace.
The split weights remain normalized and nonnegative after the actual disorder
average. Contracting the Hessian with the original SK spectral covariance
gives its overlap-moment telescope; the factor $N$ cancels exactly in the
normalized free-energy amplitude derivative.
\end{lemma}
\begin{proof}
Finite linearity of the normalized Gaussian means gives the actual probability
and moment identities. Induct the common covariance rule and telescope its
mass coefficients. The checked derivative bounds discharge all measurability
and integrability hypotheses of Gaussian coordinate and radial Stein.
\end{proof}

\begin{lemma}[Signed tagged covariance and negative field endpoints]
\lean{SpinGlass.Targets.section5Interleaved\_covariance\_increment}\leanok
\lean{SpinGlass.Targets.section5InterleavedCovarianceExpression\_le}\leanok
\lean{SpinGlass.Targets.exists\_section5Interleaved\_zero\_gap\_negative}\leanok
The signed cumulative cross path is $c_i=e\rho_{\min(i,\tau)}$, where
$e=1$ for $u\ge0$ and $e=-1$ otherwise. At each tag,
$\Delta c_i=\Delta\rho_i$ times its actual correlation, equal to $0$, $1$
or $-1$. Physical tags have $\Delta\rho_i=0$, so their sharing need not
coincide with the interpolation cutoff. For $t\ge0$ and a normalized,
nonnegative split law, the finite covariance expression is at most minus
twice the original Parisi correction times $t$.

For $r\ge2$, $m_1>0$, $q_1<q_2$, $\beta\ne0$, $0<t<1$ and $u<0$, the
actual field endpoint has a positive deficit independent of $N$ and disorder,
whenever the current trial index satisfies $1\le j\le k+1$ and
$q_{j-1}\le |u|\le q_j$ (here $k$ is the local scheme index).
This includes $q_1=0$, $j=1$ and $|u|=q_1$ when admissible.
\end{lemma}
\begin{proof}
The tagged increment identity and the existing finite covariance telescope
allow signed cross increments without changing frozen physical sharing.
Square completion uses the same signed trial matrix and its checked endpoint
correction. For the negative field endpoint, explicit original tags supply
active ordinary-shared and opposite-shared increments. Exact second-replica
reflection gives the dual strict comparison, including zero mass at the
outer opposite level. The actual zero-endpoint bound transports the strict
scalar deficit to $\eta(0)$.
\end{proof}

\begin{lemma}[Actual simultaneous mixed derivative and disorder average]
\lean{SpinGlass.Targets.hasDerivAt\_mixedConstrainedCascade\_path\_decomposition}\leanok
\lean{SpinGlass.Targets.mixedConstrainedGaussian\_path\_derivative}\leanok
\lean{SpinGlass.Targets.hasDerivAt\_mixedConstrainedGaussian\_path\_trace}\leanok
For the actual mixed cascade, independent disorder and active variance
coordinates are jointly differentiable. Along a differentiable path with
positive or locally constant zero variances, the time derivative is the
disorder direction plus the sum of original-level variance derivatives.
An anchored increment bound, affine in the disorder norm on a single time
neighborhood, justifies passage through the Gaussian expectation.
Radial Stein identifies the disorder contribution with its original spectral
Hessian trace, multiplied by $a'(w)a(w)$ for disorder amplitude $a(w)$.
Inactive zero variances contribute zero; no ordinary variance derivative
at zero is assumed.
\end{lemma}
\begin{proof}
Reuse the finite-dimensional Gaussian parameter theorem, then restrict to
the active variance face and apply the genuine chain rule. Difference
quotients give disorder measurability of the derivative. The anchored
bound is integrable under the actual Gaussian law. Proved integrability
of the disorder direction and individual heat terms permits splitting the
expectation before applying the existing radial Stein identity.
\end{proof}

\begin{lemma}[Actual mixed heat, covariance inequality and endpoint transport]
\lean{SpinGlass.Targets.mixedLevelVarianceD\_overlap}\leanok
\lean{SpinGlass.Targets.mixedConstrained\_trace\_sub\_heat\_eq\_covariance}\leanok
\lean{SpinGlass.Targets.hasDerivAt\_section5InterleavedInterpolation\_replica}\leanok
\lean{SpinGlass.Targets.deriv\_section5InterleavedInterpolation\_le}\leanok
\lean{SpinGlass.Targets.section5InterleavedInterpolation\_endpoint\_bound}\leanok
For the actual tagged construction, $1\le j\le k+1$ and
$q_{j-1}\le |u|\le q_j$, the simultaneous mixed free-energy derivative is
the signed covariance expression under the genuine split law, averaged
over the original disorder. In the exact-covariance SK setting, for
$N>0$, $0\le t\le1$ and $0<w<1$,
\[
 \eta'(w)\le -2t\,\mathrm{correction}(s,\beta).
\]
For $1\le r\le k+1$, closed-interval continuity and the actual endpoint
identity therefore imply
\[
 \Psi(t,u)\le\eta(0)-2t\,\mathrm{correction}(s,\beta)
 \le 2\psi(t)-\delta_{r,j}(t,u),
\]
where $\delta_{r,j}=2A_0(h)-V_{r,j}(0)$ is the already constructed scalar
deficit and has no system-size or disorder argument.
\end{lemma}
\begin{proof}
Uniqueness of genuine perturbation derivatives identifies packed Gaussian
means with the actual mixed normalized means. Contracting the spatial
heat directions gives correlations $0,1,-1$ and the factor $N/2$.
Unchanged outer levels preserve the same original split law as the SK
trace. The finite covariance telescope combines them after normalization
by $1/N$. Zero visited variances have zero actual velocity, by uniqueness
of the derivative of their constant paths, so the active-face indicator
can be removed without a boundary derivative. Signed square completion
and the tagged correction identity give the derivative bound. Apply the
closed-interval mean-value inequality and the checked endpoints.
\end{proof}

\begin{lemma}[Strict original-free-energy deficits in the current trial range]
\lean{SpinGlass.Targets.exists\_constrainedPhi\_interleaved\_gap\_left\_outside}\leanok
\lean{SpinGlass.Targets.exists\_constrainedPhi\_interleaved\_gap\_right\_outside}\leanok
\lean{SpinGlass.Targets.exists\_constrainedPhi\_interleaved\_gap\_negative}\leanok
The left and right outside-index scalar deficits with their stated positive
witness variances now give $\Psi(t,u)\le2\psi(t)-\delta$ for a fixed
$\delta>0$ chosen before all system sizes and SK disorders. The same holds
for the negative-overlap scalar deficits at $r\ge2$, $m_1>0$,
$q_1<q_2$, $\beta\ne0$ and $0<t<1$, within $1\le j\le k+1$.
The first physical level is also covered when $q_1>0$ and $|u|>q_1$,
under the corresponding strict negative-endpoint hypotheses.
\end{lemma}
\begin{proof}
Use the previously proved positivity of the explicit scalar deficit and
the actual endpoint bound above, without changing the deficit or its
quantifier order.
\end{proof}

\begin{lemma}[Breakpoint-safe positive outside-neighbor bounds]
\lean{SpinGlass.Targets.section5OutsideLeft\_exists\_interval}\leanok
\lean{SpinGlass.Targets.section5OutsideRight\_exists\_interval}\leanok
\lean{SpinGlass.Targets.exists\_constrainedPhi\_gap\_left\_outside}\leanok
\lean{SpinGlass.Targets.exists\_constrainedPhi\_gap\_right\_outside}\leanok
For $\beta\ne0$, $0<t<1$, strictly ordered masses and the relevant
positive physical overlap gap, the strict original-free-energy bound holds
for $0\le u<q_{r-1}$ and for $q_{r+1}<u\le q_{k+1}$, including exact
breakpoint overlaps. No trial index or positivity of auxiliary witness
variances is assumed in these final wrappers. The deficit is chosen before
system size and disorder.
\end{lemma}
\begin{proof}
Choose the first index with $u<q_j$ on the left, so that
$q_{j-1}\le u<q_j$. On the right choose the first index with $u\le q_j$,
so that $q_{j-1}<u\le q_j$. Scheme monotonicity locates the index outside
the neighboring range; the strict endpoint of the selected interval gives
the required active witness variance. Apply the transported strict bound.
\end{proof}

\begin{lemma}[Exact terminal padding and full outside trial ranges]
\lean{SpinGlass.Targets.constrainedPhi\_padOneLast}\leanok
\lean{SpinGlass.Targets.guerraPsi\_padOneLast}\leanok
\lean{SpinGlass.Targets.exists\_constrainedPhi\_gap\_right\_outside\_full}\leanok
\lean{SpinGlass.Targets.exists\_constrainedPhi\_gap\_negative\_all\_trials}\leanok
Redundant terminal mass-one padding leaves the original constrained free
energy and $\psi$ unchanged, after shifting the independent cutoff by one.
The positive outside-neighbor bound therefore covers
$q_{r+1}<u\le1$. The negative bound at $r\ge2$ covers every
$-1\le u<0$, including zero first overlap and all trial breakpoints,
under its original first-mass and first-interior-overlap gap hypotheses.
\end{lemma}
\begin{proof}
Delete the extra innermost zero-variance Gaussian step; the added correction
summand is zero. For positive strictness, the actual order obstruction needs
only $m_r<m_{j-1}$, a pair of original masses. It does not require strict
monotonicity of the padded mass sequence. The negative strict comparison
already needs only $m_1>0$ and $q_1<q_2$. Apply the proved mixed interpolation
to the padded scheme, select the trial interval by first crossing, and
transport back exactly. No minimality of the padded scheme is used.
\end{proof}

\begin{lemma}[The zero-first-overlap signed boundary]
\lean{SpinGlass.Targets.field\_eq\_zero\_of\_initial\_overlap\_zero\_min}\leanok
\lean{SpinGlass.Targets.constrainedPhi\_initial\_zero\_reflection\_of\_min}\leanok
\lean{SpinGlass.Targets.exists\_constrainedPhi\_negative\_initial\_zero\_neighbor}\leanok
At $q_1=0$, original fixed-level minimality, $\beta\ne0$, a positive
first mass and $q_2>0$ imply $h=0$. The actual first-level constrained
free energy is then invariant under $u\mapsto-u$. The previously checked
positive local-right and far-right bounds give a strictly positive
size-independent deficit on $[-q_2,0)$, under the same far-right
near-optimality smallness condition.
\end{lemma}
\begin{proof}
The actual initial scalar slope vanishes by stationarity. It is strictly
increasing and vanishes at zero by evenness, so $h=0$. The sole outer
shared variance is zero and deletes exactly. Reflect the second replica
through the independent Gaussian cascade using SK disorder evenness.
This is an identity for the original free energy, not a replacement field
model. Apply the existing positive estimates at $-u$.
\end{proof}

\begin{proposition}[Proposition 5.7 for reduced near-minimizers in the SK proof]
\lean{SpinGlass.Targets.exists\_section5\_common\_accuracy}\leanok
\lean{SpinGlass.Targets.exists\_constrainedPhi\_gap\_outside\_of\_reduced\_min}\leanok
\lean{SpinGlass.Targets.talagrand\_proposition\_5\_7}\leanok
For $\beta\ne0$ and $t_0<1$, there is an $\varepsilon>0$, depending
only on $\beta,t_0$, with the following property. Let $s$ be a fixed-level
minimizer within $\varepsilon$ of the infimum, with strictly increasing
masses and strictly increasing interior overlaps. Boundary overlaps
$q_1=0$ and $q_{k+1}=1$ are allowed. For every $1\le r\le k+1$,
$0\le t\le t_0$ and $-1\le u\le1$ such that
$u<q_{r-1}$ or $u>q_{r+1}$, there is a $\delta>0$ satisfying
\[
  \Psi(t,u)\le 2\psi(t)-\delta
\]
for every $N>0$, every SK disorder, and every attainable such overlap.
The choice of $\delta$ precedes $N$ and the disorder.
\end{proposition}
\begin{proof}
Continuity of the sixth root and square root at zero supplies one positive
accuracy satisfying both Section~5 smallness conditions before the scheme
is chosen. For $t>0$, use the full positive outside-neighbor comparisons
and negative interleaved bounds. At the first physical level, Proposition~5.4
includes the endpoint $u=-q_1$; the strict mixed bound covers $u<-q_1$
when $q_1>0$. When $q_1=0$, use the exact reflection above, the neighboring
negative bound, and the full positive outside bound. At $t=0$, the previously
proved quadratic deficit applies and the outside condition ensures $u\ne q_r$.
\end{proof}

\begin{lemma}[Continuous scalar witnesses and compact subregions]
\lean{SpinGlass.Targets.continuousOn\_section5InterleavedScalarV\_trial}\leanok
\lean{SpinGlass.Targets.section5InterleavedScalarV\_time\_zero}\leanok
\lean{SpinGlass.Targets.exists\_uniform\_positive\_of\_compact\_witnesses}\leanok
\lean{SpinGlass.Targets.exists\_constrainedPhi\_compact\_left\_gap}\leanok
\lean{SpinGlass.Targets.exists\_constrainedPhi\_compact\_right\_gap}\leanok
\lean{SpinGlass.Targets.exists\_uniform\_constrainedPhi\_left\_outside\_trial}\leanok
For each fixed lambda, the actual mixed scalar comparison is continuous in
$(t,u)$ on its closed admissible trial strip, including the sign boundary.
At time zero the full-lambda family equals the physical stationary scalar
family; consequently its quadratic lambda gain is available as a continuous
witness, including after exact terminal padding using original minimality.

More generally, let $S$ be compact and let $(D_i)_i$ be continuous scalar
deficits on $S$. If for every $z\in S$ some $D_i(z)>0$, then there is one
$\delta>0$ such that every $z\in S$ has some $D_i(z)\ge\delta$.
For the actual comparison deficits this gives
$\Psi(t,u)\le 2\psi(t)-\delta$, uniformly before $N$ and the disorder.
The checked concrete applications cover compact left/right neighboring
regions satisfying the far-gap smallness hypotheses and compact left
outside-trial regions $q_{j-1}\le u<q_j$, $j<r$, including time zero,
under strict masses, the physical overlap gap and stationarity.
\end{lemma}
\begin{proof}
Specialize the existing good scalar families to time and overlap. Zero
variances delete exactly and identify the two sign charts at zero.
For compactness, freeze a strictly positive witness at each point; its
half-value persists on an open neighborhood. A finite subcover and the
minimum of these positive half-values give $\delta$. Witness selection
need not be continuous. The actual arbitrary-lambda endpoint inequality
transports the scalar bound to the constrained free energy.
\end{proof}

The compact uniform quadratic bound of Theorem~2.4 is now checked by the
public Lean theorem cited below. Its finite assembly joins adjacent trial
indices at shared breakpoints and includes the signed, first-level,
near-terminal and terminal domains. No continuity is asserted for the
changing finite-size overlap constraint.

\begin{lemma}[Adjacent inserted-overlap boundary]
\lean{SpinGlass.Targets.section5Rho\_adjacent\_boundary}\leanok
\lean{SpinGlass.Targets.section5TaggedMass\_adjacent\_boundary\_of\_ne}\leanok
\lean{SpinGlass.Targets.section5SortedTagList\_adjacent\_boundary\_erase}\leanok
\lean{SpinGlass.Targets.section5InterleavedScalarV\_adjacent\_boundary}\leanok
\lean{SpinGlass.Targets.section5InterleavedLambdaDeficit\_adjacent\_boundary}\leanok
At the shared breakpoint $u=q_j$, inserting $u$ at level $j$ or at level
$j+1$ gives the same cumulative overlap sequence:
\[
  \rho_{j}(p;q_j)=\rho_{j+1}(p;q_j)\qquad(p\ge0).
\]
The changed interpolation tag has zero variance at the breakpoint. Stable
tagged sorting and exact zero-variance deletion therefore identify the full
scalar recursion, and the arbitrary-lambda deficit transports across the
boundary.
\end{lemma}

The compact right outside-trial region is handled symmetrically, with a
half-open lower overlap endpoint so that the interpolating variance is
positive. Mathlib's stable insertion-sort permutation lemma, scalar-list
zero-variance filter, reverse-index bridge, terminal padding, finite compact
cover bookkeeping, and positive-time negative-overlap compact adapter are
checked. These pieces are assembled into the single quadratic estimate by
\lean{SpinGlass.Targets.talagrand\_theorem\_2\_4}\leanok.

\begin{lemma}[Closed-interval concavity and the supporting line (5.34)]
\lean{SpinGlass.Targets.concaveOn\_section4U}\leanok
\lean{SpinGlass.Targets.section4U\_sub\_ge\_supportingLine}\leanok
\lean{SpinGlass.Targets.section4U\_interpolation\_slope\_lower\_bound}\leanok
For $m_{r-1}<1$, the actual $U$ is concave on $[0,a]$. For $v\in[0,a]$
and $t\in[0,1]$,
\[
 U(tv)\ge U(v)-(1-t)vQ(m_{r-1},tv).
\]
Consequently, writing $v=\beta^2(q_r-u)$ and
$D=U(tv)-t\beta^2(q_r^2-u^2)/2$,
\[
 D\ge2f(u)+(1-t)\frac{\beta^2}2(q_r-u)^2
       -(1-t)v\bigl(Q(m_{r-1},tv)-u\bigr).
\]
Both time and variance endpoints, including a zero variance gap, are allowed.
\end{lemma}
\begin{proof}
Reuse Mathlib's derivative criterion for concavity with $U'=Q$ and
$Q'=-R\le0$ on the interior, and the proved closed-interval continuity.
The inward derivative identifies the supporting slope at either endpoint.
Apply the supporting line at $tv$ to $v$ and expand the SK correction.
\end{proof}

\begin{lemma}[Actual constrained free-energy gains]
\lean{SpinGlass.Targets.constrainedPhi\_le\_two\_guerraPsi\_sub\_factor\_sq}\leanok
\lean{SpinGlass.Targets.constrainedPhi\_le\_two\_guerraPsi\_mass\_taylor}\leanok
\lean{SpinGlass.Targets.constrainedPhi\_le\_guerraPsi\_right\_lambda\_gain}\leanok
For an attainable overlap in the left interval, set $v=\beta^2(q_r-u)$.
The actual constrained free energy satisfies
\[
 \Psi(t,u)\le2\psi(t)-\tfrac12\bigl(Q(m_{r-1},tv)-u\bigr)^2.
\]
For $m_{r-1}/2\le m\le m_r$, with $D$ as above, it also satisfies
\[
 \Psi(t,u)\le2\psi(t)+D(m-m_{r-1})+2C(\beta)(m-m_{r-1})^2.
\]
On the right interval the analogous lambda bound holds with the actual
right-transform derivative in place of $Q$; its dual scalar identification
is not assumed.
\end{lemma}
\begin{proof}
Combine the actual second-interpolation endpoint inequalities with the
time-zero lambda bounds. At baseline mass the correction is exactly that
in $2\psi(t)$. For the mass bound use $V(0,m,tv)=2T(tv,m)$ and the full
two-sided mass Taylor estimate, including decreases from positive baseline.
\end{proof}

\begin{lemma}[Positive-baseline far-left strict improvement]
\lean{SpinGlass.Targets.constrainedPhi\_lt\_two\_guerraPsi\_of\_left\_gap}\leanok
\lean{SpinGlass.Targets.exists\_constrainedPhi\_left\_gap\_uniform\_in\_size}\leanok
Under the actual optimality hypotheses of Proposition~4.6, suppose
$0<m_{r-1}<m_r$ and
\[
 2L(\beta)\sqrt\varepsilon<(1-t)\frac{\beta^2}2(q_r-u)^2.
\]
Then $\Psi(t,u)<2\psi(t)$. More precisely a positive deficit can be chosen
before $N$ and the SK disorder, for the fixed scheme, $t$ and $u$.
\end{lemma}
\begin{proof}
If $Q(m_{r-1},tv)\ne u$, the lambda deficit is positive. Otherwise the
supporting-line bound and Proposition~4.6 give $D>0$. Choose
\[
 \delta=\min\left\{\frac{m_{r-1}}2,\frac{D}{4(C(\beta)+1)}\right\}>0.
\]
The admissible mass $m=m_{r-1}-\delta$ gives a deficit at least $D\delta/2$.
Both deficits are independent of $N$. This is the positive-baseline
argument of Proposition~5.5, not the initial mass-zero argument, nor the
compactness and local quadratic estimates required by Theorem~2.4.
\end{proof}
\lean{SpinGlass.Targets.constrainedPhi\_lt\_two\_guerraPsi\_of\_far\_left}\leanok
\lean{SpinGlass.Targets.exists\_constrainedPhi\_far\_left\_uniform\_in\_size}\leanok
The paper's uniform-smallness form is also checked: choose
$L_3(\beta)=4L(\beta)$. For $\beta\ne0$, $L_1>0$, $0\le t\le t_0<1$,
\[
 L_1(q_r-u)\ge1-t_0,\qquad
 L_3(\beta)\sqrt\varepsilon\le
 (1-t_0)\frac{\beta^2}2\left(\frac{1-t_0}{L_1}\right)^2
\]
imply the preceding strict bound under the same optimality and positive-mass
hypotheses. The constant $L_3$ is independent of depth, field and gap.

\begin{lemma}[Actual normalized constrained replicas at every split]
\lean{SpinGlass.Targets.sum\_constrainedCascadeGibbs}\leanok
\lean{SpinGlass.Targets.sum\_constrainedCascadeReplica}\leanok
\lean{SpinGlass.Targets.constrainedCascadeReplica\_disorder\_product}\leanok
\lean{SpinGlass.Targets.constrainedCascadeReplica\_spatial\_product}\leanok
The terminal Gibbs coordinates transported through each inner depth are
nonnegative normalized probabilities. Form their product at a given depth,
then transport that product through the actual remaining outer levels.
These split-level replica weights are again normalized and represent
products of the genuine inner disorder and spatial first derivatives.
\end{lemma}
\begin{proof}
Normalized positive tilted means preserve constants and nonnegativity.
Finite linearity gives both normalization and moment identities. Restarting
at depth $l$ uses the actual potential there and the shifted cutoff $d-l$.
SK spectral contraction gives $N$ times the four-overlap covariance;
independent and signed shared field contractions give the corresponding
diagonal and cross-overlap sums. Products are formed before outer averaging,
not from two already averaged probabilities.
\end{proof}

\begin{lemma}[Actual full replica Hessian and averaged SK trace]
\lean{SpinGlass.Targets.constrainedPairFieldCascadeSecond\_eq\_replica}\leanok
\lean{SpinGlass.Targets.constrainedPairCascadeSpatialSecond\_eq\_replica}\leanok
\lean{SpinGlass.Targets.sum\_averagedConstrainedCascadeReplica}\leanok
\lean{SpinGlass.Targets.integral\_constrainedPairFieldCascadeSecond\_SK\_trace}\leanok
Write $P_j$ for the actual depth-$j$ Gibbs law and $\mu_{l,j}$ for the
actual split-$l$ two-replica law transported to total depth $j$. For terminal
direction observables $f,g$, the genuine nested Hessian is
\[
 \langle fg\rangle_{P_j}-\mu_{0,j}(f\otimes g)
 +\sum_{l<j}m_l\bigl(\mu_{l,j}(f\otimes g)-\mu_{l+1,j}(f\otimes g)\bigr).
\]
The same formula holds for spatial direction observables. Contracting the
disorder directions with the actual SK spectral covariance, with kernel
$K(p,q)=\beta^2\sum_{a,b}R_{ab}(p,q)^2/2$, gives the outer-averaged trace
\[
 N\left(\beta^2(1+u^2)-\E\mu_{0,j}(K)
 +\sum_{l<j}m_l\bigl(\E\mu_{l,j}(K)-\E\mu_{l+1,j}(K)\bigr)\right).
\]
\end{lemma}
\begin{proof}
For an independent level the intermediate covariance products of its two
equal-mass tilts cancel pointwise, giving one covariance increment.
Expand the genuine Hessian recursion using bounded-observable linearity and
the exact shifted-level transport. The terminal covariance supplies the first
two terms. Joint disorder/field measurability and the unit bound justify
every outer expectation and finite-moment interchange. The SK contraction
and constrained diagonal give the displayed trace, including its factor $N$.
All fixed masses and variances may vanish. Adjacent-mass telescoping retains
the initial and final coefficients rather than assuming endpoint masses.
\end{proof}
\begin{lemma}[Actual original-level heat generators under the split law]
\lean{SpinGlass.Targets.constrainedLevelVarianceD\_independent\_eq\_replica}\leanok
\lean{SpinGlass.Targets.constrainedLevelVarianceD\_shared\_eq\_replica}\leanok
Each actual original-level heat generator, after every remaining outer level,
is one half the coordinate sum of the transported inner replica Hessian plus
the original mass times the transported inner squared first derivative.
Independent levels use the two separate spin directions; shared levels use
the common direction. Every term uses the same actual split law as above.
\end{lemma}
\begin{proof}
Exact transport composition increases only the remaining outer depth, retaining
each product's original split index. Identify packed independent Gaussian means
with nested equal-mass means by uniqueness of their genuine bounded-perturbation
derivatives. Then apply the checked original-level heat-generator identities
and finite linearity. The algebra includes zero variances but does not claim
a derivative at zero variance. The next result contracts the field directions,
justifies the disorder average and combines their physical coefficients.
\end{proof}

\begin{lemma}[Actual covariance identity and Section 5 endpoint inequalities]
\lean{SpinGlass.Targets.integral\_constrainedLevelVarianceD\_independent\_overlap}\leanok
\lean{SpinGlass.Targets.integral\_constrainedLevelVarianceD\_shared\_overlap}\leanok
\lean{SpinGlass.Targets.averagedReplicaPressureDerivative\_eq\_covariance}\leanok
\lean{SpinGlass.Targets.section5Interpolation\_endpoint\_bound}\leanok
\lean{SpinGlass.Targets.section5RightInterpolation\_endpoint\_bound}\leanok
The actual field heat terms contract to the independent kernel
$R_{11}+R_{22}$ or the shared kernel $\sum_{a,b}R_{ab}$, with factor $N/2$.
After disorder averaging and summation over levels, both physical Section~5
derivatives equal the covariance expression $\mathcal B$ above. Consequently,
for the allowed left mass interval $m_{r-1}/2\le m\le m_r$,
\[
 \eta(1)\le\eta(0)-t\left(2\mathcal C_s+
 (m-m_{r-1})\frac{\beta^2}{2}(q_r^2-u^2)\right),
\]
where $\mathcal C_s$ is the original Parisi correction and
$\eta(1)$ is the actual constrained pressure. This is (5.9).
For the dual right interval, $m_{r-1}\le m\le2m_r$, the added correction is
$(m-m_r)\beta^2(u^2-q_r^2)/2$.
\end{lemma}
\begin{proof}
The constrained diagonals are $2$ and $2(1+u)$. Actual normalized split
weights justify all finite moment interchanges. The remaining field
increments telescope to the trial covariance at the original split, while
summation by parts retains and then cancels the endpoint mass coefficients.
Reverse indices to the paper's convention and use the checked square
completion. At zero variance the actual velocity is zero. Finally use
closed-interval continuity and the mean-value theorem; no derivative at an
interpolation endpoint is required. These are the exact-covariance SK,
attainable nonnegative-overlap constructions, not the general signed case.
\end{proof}

\begin{lemma}[One-level heat generators for actual constrained cascades]
\lean{SpinGlass.Targets.hasDerivAt\_sharedStepPi\_constrained\_variance}\leanok
\lean{SpinGlass.Targets.hasDerivAt\_independentStepPi\_constrained\_variance}\leanok
For a fixed inner constrained cascade $F$, an added Gaussian level of actual
mass $m$ and variance $v>0$ has derivative
\[
 \tfrac12\mathbb E_W\sum_i\bigl(D_i^2F+m(D_iF)^2\bigr).
\]
The shared case uses common replica-field directions; the independent case
uses the $2N$ separate directions. Zero added mass is included.
\end{lemma}
\begin{proof}
Prove linearity of the propagated spatial directions, apply the existing
parameter chain rule to the Gaussian amplitude, and use normalized Gaussian
Stein. The actual cascade's previously checked growth and Hessian bounds
discharge all analytic assumptions. Mathlib's measure-preserving product
reindexing packs the independent fields into $2N$ coordinates.
\end{proof}
\begin{lemma}[Actual averaged partial derivatives of the full paired cascade]
\lean{SpinGlass.Targets.hasDerivAt\_constrainedFieldCascade\_variance}\leanok
\lean{SpinGlass.Targets.hasDerivAt\_constrainedDisorderPressure}\leanok
\lean{SpinGlass.Targets.hasDerivAt\_constrainedDisorderPressure\_variance}\leanok
Fix all paired masses and field variances. The disorder-only pressure
$P(t)=N^{-1}\mathbb E F(\sqrt t H)$ of the actual nested constrained cascade has
\[
 P'(t)=\frac1{2N}\mathbb E\sum_i\tau_i D_{w_i}^2F(\sqrt t H),\qquad t>0.
\]
It is continuous at $t=0$. Separately, changing one original field variance
to $v>0$ differentiates its actual heat seed through every unchanged outer
level. The resulting derivative is an explicit nested normalized mean; its
bound is $NK$ for an independent level and $2NK$ for a shared level, where
\[
 K=2+4\sum_{i<\ell}|m_i|+|m_\ell|.
\]
This bound has no additional outer-depth factor. The derivative also passes
through the actual Gaussian disorder average and site normalization.
\end{lemma}
\begin{proof}
For disorder, prove finite linearity of the actual propagated direction, then
apply scaled coordinate Stein. The mixed-Hessian bound and Gaussian norm
integrability justify differentiation under the expectation. Compose with
$t\mapsto\sqrt t$ for the coefficient $1/(2N)$ and endpoint continuity.
For a field variance, use the existing heat generator at the chosen level and
the normalized parameter rule at every outer level. Measurable forward
difference quotients give disorder measurability of the actual variance
derivative; its uniform bound justifies the outer Gaussian differentiation.
\end{proof}
The full replica-covariance identification and endpoint-safe integration for
the physical positive-overlap paths are now proved above.

\begin{lemma}[Degenerate interpolation variances need no endpoint derivative]
\lean{SpinGlass.Targets.section5InterpolationVariance\_pos\_or\_eq\_zero}\leanok
\lean{SpinGlass.Targets.section5RightInterpolationVariance\_pos\_or\_eq\_zero}\leanok
Before interpolation time one, each actual field variance in either neighbor
construction is positive or identically zero. Coincident overlaps are allowed.
\end{lemma}
\begin{proof}
Each variance is $A+(1-w)B$ with $A,B\ge0$. If this vanishes at $w<1$, then
$A=B=0$. Such a coordinate is constant and can be omitted from the eventual
chain rule; no two-sided variance derivative at zero is required.
\end{proof}

\begin{lemma}[Moving coefficients and actual second-interpolation continuity]
\lean{SpinGlass.Targets.hasDerivAt\_section5InterpolationAmplitude}\leanok
\lean{SpinGlass.Targets.continuousOn\_constrainedPairFieldCascade\_paths}\leanok
\lean{SpinGlass.Targets.continuousOn\_section5Interpolation}\leanok
\lean{SpinGlass.Targets.continuousOn\_section5RightInterpolation}\leanok
For each actual field variance $v(w)=A+(1-w)B$, its amplitude derivative
before time one is $-B/(2\sqrt{v(w)})$. A zero face has $A=B=0$, so this
totalized expression agrees with the derivative of a constant.
The disorder coefficient $\sqrt{wt}$ is similarly differentiated for $w>0$,
including $t=0$ by its constant branch.

The actual constrained cascade is continuous along simultaneous continuous
disorder, field and nonnegative variance paths on compact parameter sets.
After Gaussian disorder averaging, both physical second interpolations are
continuous on the entire interval $[0,1]$.
\end{lemma}
\begin{proof}
Generalize the parameter type of the existing Gaussian continuity theorem,
without changing its proof. Reuse it at variance 1 with moving coefficients
in the input, using compact-image growth bounds and the existing independent/
shared packing identities. RSAT's finite-state log partition supplies the
actual terminal continuity. The disorder Lipschitz estimate gives an affine
Gaussian-norm dominating function for the outer expectation.
\end{proof}
These results do not assert the simultaneous derivative or its covariance bound.

\begin{lemma}[Joint differentiation of the actual moving cascade]
\lean{SpinGlass.Targets.contDiff\_constrainedPairFieldBase\_joint}\leanok
\lean{SpinGlass.Targets.differentiableAt\_constrainedFieldCascade\_joint}\leanok
\lean{SpinGlass.Targets.hasDerivAt\_constrainedFieldCascade\_path\_fderiv}\leanok
\lean{SpinGlass.Targets.differentiableAt\_section5Interpolation\_integrand}\leanok
\lean{SpinGlass.Targets.differentiableAt\_section5RightInterpolation\_integrand}\leanok
Let the disorder and all visited variances follow differentiable paths near
$w$. Assume the variances are nonnegative nearby and that each visited
coordinate is positive at $w$ or locally identically zero. The actual
constrained cascade is jointly differentiable in the path parameter and both
replica fields. In particular, both physical second-interpolation integrands
are differentiable for $0<w<1$, at fixed disorder, including zero masses and
coincident overlaps.
\end{lemma}
\begin{proof}
The absorbed finite-state log partition gives joint terminal smoothness.
Telescope the individual variance comparisons and combine them with the
disorder and spatial Lipschitz bounds to obtain an anchored parameter bound,
uniform in the replica fields. At each Gaussian step differentiate the moving
shift directly. Its difference bound is affine in the Gaussian $\ell^1$ norm;
the log-Laplace branch has an integrable exponential-growth bound. Mathlib's
dominated differentiation therefore applies to the full joint input. Handle
mass zero as an ordinary expectation and use the actual independent/shared
packing for induction. The positive-or-constant-zero alternative verifies the
physical paths' hypotheses.
\end{proof}
\begin{lemma}[The actual Gaussian-averaged simultaneous derivative]
\lean{SpinGlass.Targets.constrainedPairGaussian\_path\_derivative}\leanok
\lean{SpinGlass.Targets.hasDerivAt\_section5Interpolation}\leanok
\lean{SpinGlass.Targets.hasDerivAt\_section5RightInterpolation}\leanok
The simultaneous path derivative passes through the actual outer Gaussian
average. Thus both physical second interpolations are differentiable on
$(0,1)$, retaining the exact normalized integral of the joint derivative.
\end{lemma}
\begin{proof}
Choose the amplitude and finite-variance bounds before fixing the disorder:
on one common time neighborhood the actual difference is bounded by
$(A+B\|U\|)|z-w|$. This is Gaussian integrable. Measurable difference quotients
give measurability of the actual derivative; clipping variances to their
nonnegative parts is used only in the auxiliary quotients and agrees with the
original path nearby. Mathlib's anchored dominated differentiation proves
integrability and interchange. Substitute the physical paths and their
distinct independent/shared cutoffs, preserving the factor $1/N$.
\end{proof}
\begin{lemma}[Explicit averaged trace and level-heat contributions]
\lean{SpinGlass.Targets.hasDerivAt\_constrainedFieldCascade\_path\_decomposition}\leanok
\lean{SpinGlass.Targets.hasDerivAt\_constrainedPairGaussian\_path\_trace}\leanok
\lean{SpinGlass.Targets.hasDerivAt\_section5Interpolation\_trace}\leanok
\lean{SpinGlass.Targets.hasDerivAt\_section5RightInterpolation\_trace}\leanok
For the actual cascade $F_j$ with disorder amplitude $a(w)$ and level
variances $v_l(w)$, the Gaussian-averaged derivative is
\[
 a'(w)a(w)\,\mathbb E\!\sum_i\tau_i
   D^2_{e_i,e_i}F_j
   +\sum_{l<j}v_l'(w)\,\mathbb E H_l.
\]
Here $H_l$ is the proved derivative in original level $l$ at positive variance;
locally constant zero coordinates contribute zero. For both physical
interpolations, site normalization gives disorder coefficient $t/(2N)$.
\end{lemma}
\begin{proof}
Prove finite-dimensional joint regularity by the same dominated Gaussian
induction. Mask the zero coordinates and identify the Fr\'echet derivative's
basis values with the existing disorder and individual variance derivatives
by uniqueness. Finite linearity gives the actual path sum. Each term is
Gaussian integrable; split the expectation and apply the existing radial
Stein identity. The physical coefficient satisfies $a'a=t/2$, including
$t=0$ by its constant branch. The heat normalization is already contained in $H_l$.
\end{proof}
The actual replica-overlap identification, covariance inequality and endpoint
transport for both physical positive-overlap paths are now proved above.

\begin{lemma}[Right-interval endpoints, baseline and scalar gain]
\lean{SpinGlass.Targets.section5RightInterpolation\_one}\leanok
\lean{SpinGlass.Targets.section5RightInterpolation\_zero\_le}\leanok
\lean{SpinGlass.Targets.section5RightInterpolation\_zero\_lambda\_gain}\leanok
For $q_r\le u\le q_{r+1}$, the dual construction has an inserted shared level
of paired mass $m/2$ and split variance $t\beta^2(u-q_r)$. Its time-one value
is the original constrained pressure, and its time-zero value has the scalar
comparison analogous to (5.17). At the baseline $m=m_r$,
\[
 V_{\rm right}(0,m_r,v)=2A_0(h),\qquad
 \eta_{\rm right}(0)\le2(\log2+A_0(h))
       -\tfrac12(V_{\rm right}'(0)-u)^2.
\]
Both variance endpoints and the last overlap interval are included.
\end{lemma}
\begin{proof}
The inserted frozen variance is zero. Move the sharing cutoff across this
redundant level and reuse the existing time-one equality. Reuse site
tensorization for time zero, the scalar semigroup for the baseline, and the
sharp lambda-curvature invariant for the explicit optimized gain.
\end{proof}
The gain is now transported to the actual constrained free energy above.
Identification of the lambda derivative with the dual $U'$ remains open.

\begin{lemma}[Dual correction and full right scalar comparison]
\lean{SpinGlass.Targets.section5RightCorrection\_eq}\leanok
\lean{SpinGlass.Targets.hasDerivAt\_section4RightCascade\_split}\leanok
\lean{SpinGlass.Targets.antitoneOn\_section4RightT}\leanok
The actual dual paired correction, with shared cutoff $r+1$, is
\[
 2\operatorname{ParisiCorrection}
 +(m-m_r)\frac{\beta^2}{2}(u^2-q_r^2).
\]
For $0\le m\le m_r$, the full actual $T_{\rm right}(v,m)$ is continuous and
nonincreasing on the closed split interval, with
$T_{\rm right}(0,m)=A_0(h)$ and hence $T_{\rm right}(v,m)\le A_0(h)$.
The actual two-step derivative has coefficient $(m-m_r)/2$ times the
normalized tilted squared first derivative.
\end{lemma}
\begin{proof}
Reuse the generic shared/independent correction sum, changing only the two
pieces of $[q_r,q_{r+1}]$. Reflect (4.11) by $v\mapsto a-v$ with inner mass
$m_r$ and outer mass $m$, then propagate the closed-interval comparison through
order-preserving outer scalar steps. The last overlap interval is included.
\end{proof}
\begin{lemma}[Full right-recursion derivative and closed-interval bound]
\lean{SpinGlass.Targets.hasDerivAt\_section4RightT\_variance}\leanok
\lean{SpinGlass.Targets.section4RightT\_variance\_dist\_le}\leanok
For $0\le m\le m_r$, the full actual dual recursion has interior derivative
in $[(m-m_r)/2,0]$. Its Lipschitz constant on the closed split interval is
$(m_r-m)/2$, independent of the number of outer levels. The last overlap
interval and zero masses are included.
\end{lemma}
\begin{proof}
Propagate the reflected two-step derivative through exactly $r$ normalized
outer means. Their bound is preserved, and the closed-interval comparison
gives the sign. Apply the mean-value theorem with endpoint continuity.
\end{proof}
Dual mixed mass/variance identities and quantitative stationarity remain open.

\begin{lemma}[Lemma 5.9 and the optimized scalar endpoint]
\lean{SpinGlass.Targets.talagrand\_lemma\_5\_9}\leanok
\lean{SpinGlass.Targets.section5Interpolation\_zero\_lambda\_gain}\leanok
For the actual Section~5 scalar recursion at the original inserted mass,
\[
 0\le \partial_\lambda^2 V\le1,
 \qquad
 \exists\lambda:\quad V(\lambda)-\lambda u
 \le V(0)-\tfrac12(V'(0)-u)^2.
\]
The curvature constant is independent of the number of levels and all parameters.
Combining the explicit choice $\lambda=u-V'(0)$ with (5.17) and (5.19) gives
\[
 \eta(0)\le2(\log2+A_0(h))-\tfrac12(V'(0)-u)^2
\]
for attainable overlaps in the constructed left interval.
\end{lemma}
\begin{proof}
At the terminal, the lambda second derivative is $E=1-D^2$.
A normalized tilt of mass $m\in[0,1]$ sends $D$ to $\mathbb E_WD$ and
$E$ to $\mathbb E_WE+m\operatorname{Var}_W(D)$. Thus
$0\le E\le1-D^2$ survives every independent and shared transform.
RSAT supplies all analytic regularity. Mathlib convexity supplies the tangent
quadratic bound, followed by the explicit lambda choice and the checked endpoint.
\end{proof}
This does not identify $V'(0)$ with $U'$ or transport the gain to $\eta(1)$.

\begin{lemma}[Conditional finite-overlap deduction of Proposition 2.3]
\lean{SpinGlass.Targets.talagrand\_proposition\_2\_3\_of\_quadratic\_bound}\leanok
\lean{SpinGlass.Targets.guerraPhi\_uniform\_of\_quadratic\_bound}\leanok
Fix a scheme with positive first mass. Suppose the Theorem~2.4 quadratic
bound holds eventually in $N$, with one $K>0$ uniform in the replica level
and $0<t<t_0<1$. For $\eta>0$, Proposition~2.5 and a finite-overlap union give
\[
 \mu_r\!\left((R-q_r)^2\ge2K(\psi-\varphi_N)+\eta\right)
 \le (N+1)C\exp(-N/C),
\]
where $C$ is the existing event constant for deficit $\eta/K$.
Consequently the per-level and mass-weighted tails are eventually at most
$\eta$, uniformly in $t$, and $\varphi_N\to\psi$ uniformly on $[0,t_0]$.
\end{lemma}
\begin{proof}
Ising overlaps have at most $N+1$ distinct values. The displayed threshold
converts the bound relative to $2\psi$ into a pressure deficit $\eta/K$
relative to $2\varphi_N$. Apply Proposition~2.5, finite additivity of the
actual replica functional, and exponential decay. Reuse the proved
mass-weighted conversion and endpoint-safe convergence argument.
\end{proof}
This certifies the deduction only: the uniform quadratic estimate remains unproved.

\begin{lemma}[Exact leading-zero-mass reduction]
\lean{SpinGlass.Targets.exists\_positive\_first\_mass\_reduction}\leanok
Every admissible scheme has an equivalent scheme of lower or equal level with
positive first mass. The Parisi functional, $\psi(t)$, and the actual $\varphi_N(t)$
for $0\le t\le1$ are preserved. Fixed-level minimality transfers exactly.
Consequently the original Theorem~2.2 quantifiers follow from a uniform quadratic
bound proved only for positive-first-mass schemes.
\end{lemma}
\begin{proof}
If the first mass vanishes, merge the two outer Gaussian expectations by the
scalar and N-site semigroup laws. Pad every smaller competitor with a redundant
zero-variance first interval to transfer fixed-level minimality. Induct on the
number of levels; at local level zero the first mass is already one.
\end{proof}
This reduction does not assert strictness of every overlap increment in (2.19).

\begin{lemma}[Exact equal-mass compression]
\lean{SpinGlass.Targets.exists\_strict\_mass\_reduction}\leanok
Every scheme has an exactly equivalent lower-or-equal-level scheme whose
entire mass sequence is strictly increasing. The functional, $\psi$, actual
$\varphi_N$, and fixed-level minimality are preserved as above.
\end{lemma}
\begin{proof}
The arbitrary-mass N-site Gaussian semigroup merges adjacent steps of equal
mass. The scalar endpoint follows by one-site tensorization, and the two
correction terms telescope. Iterate the exact deletion, using zero-variance
padding to transfer minimality at each stage.
\end{proof}
This completes the mass-strictness part of (2.19). The combined reduction
proved above also removes repeated interior overlaps while preserving strict
masses. Thus the convergence theorem requires the quadratic bound only for
these reduced schemes. The baseline mass at $r=1$ is
nevertheless $m_0=0$. Both its scalar derivative and the uniform nested
mass-variation estimates are covered by the checked results above.

The analytic work through Proposition~2.3, Theorems~2.2 and~2.4, the exact scheme
reduction and the final Parisi-formula deduction is checked. The import graph now
separates the foundational interpolation in \texttt{Targets/TalagrandCore.lean}
from the canonical final declarations in \texttt{Targets/TalagrandFinal.lean};
\texttt{Targets/Talagrand.lean} remains the public façade. The detailed record is
maintained in \texttt{docs/ROADMAP.md}.

\begin{thebibliography}{9}
\bibitem{Gue03} F.~Guerra, Broken replica symmetry bounds in the mean field spin glass
model, \emph{Comm. Math. Phys.} 233 (2003), 1--12.
\bibitem{GT02} F.~Guerra, F.~L.~Toninelli, The thermodynamic limit in mean field spin glass
models, \emph{Comm. Math. Phys.} 230 (2002), 71--79.
\bibitem{Tal06} M.~Talagrand, The Parisi formula, \emph{Ann. of Math.} 163 (2006), 221--263.
\bibitem{TalI} M.~Talagrand, \emph{Mean Field Models for Spin Glasses}, Vol.~I, Springer 2011.
\bibitem{TalII} M.~Talagrand, \emph{Mean Field Models for Spin Glasses}, Vol.~II, Springer 2011.
\bibitem{Pan13} D.~Panchenko, \emph{The Sherrington--Kirkpatrick Model}, Springer 2013.
\end{thebibliography}

\end{document}
